%% 第26回編集委員会（物理第4巻）：組 電W5
%% p4-08.tex  電-8　電気が生まれる：電磁誘導とインダクタンス（★1〜5・3段）
%% 中身は計画書 tools/phys4-plan.md の B「電-8」。公式の持ち主は A（A13・A45〜A50・A73〜A75）。
%%   ほかの組の章は、計画書 B の節の番号で指した（電-2 ⑨、電-3 ④⑦、電-4 ③⑤⑦、電-6 ③⑤⑥⑦⑧、電-7 ③⑤⑥⑧、
%%   電-9 ②④⑦、電-10 ⑦、電-12 ③⑤⑥、電-14 ③⑨、電-15 ⑩）。
%% 数値はすべて Python で計算した：checks/物理第4巻/電-8/calc_p4_08.py（出力 calc_p4_08_out.txt）、
%%   code_p4_08.py（⑪の囲みのプログラム）、kanji_p4_08_09.py（常用漢字の確かめ）。
%% 前の巻・数学の本との約束（計画書 A-6・C-6-2）
%%   ・P45：コイルのきまり V_L＝L di/dt（数学 6-4⑦ が「部品のきまり」として使った式）→ ⑥（ph:電-8:L）
%%   ・数学 6-1:359・6-4:580・6-6:747 の L・½Li²「なぜかは物理シリーズ『電磁気学』で」→ ⑥⑧
%%   ・8-7:973「起電力は閉じた道の線積分（物理 電-8、8-14）」→ ③⑨
%%   ・8-8:1331「磁束 Φ＝BA cosθ、発電機（物理 電-8）」→ ③⑤（面積は S と書き、8-8 の A と一言）
%%   ・8-11 S6（ファラデーの法則は道を選ぶ）→ ③⑨。8-14 の（3）→ ⑨
%%   ・電-14 の計画が指す「電-8 例題2」（棒 0.5 m・2 m/s・0.1 T）→ 例題2 の数をそろえた
%%   ・電-9 の計画が指す「電-8 例題6」（リレーを切ると 60 V）→ 例題6 の数をそろえた
%% 記号（この章）：磁場（磁束密度）B、磁束 Φ（Wb）、誘導起電力 \mathcal E（高校の教科書の V。Y-3）、
%%   巻数 N、1 m あたりの巻数 n、面積 S（数学 8-8 は A）、長さ l、棒の速さ v（この章では瞬時の電圧を v と書かない）、
%%   電流 I（瞬時の電流も i と書く所は⑥⑦⑧）、抵抗 R、自己インダクタンス L、相互インダクタンス M
%%   （電-7 の磁化 M とは章が別。カードに一言）、コイルの電圧 V_L（数学 6-4 と同じ）、角速度 ω、周波数 f、
%%   磁気エネルギー U・密度 u、導電率 σ、透磁率 μ・比透磁率 μ_r、表皮深さ δ、板の厚さ d、電場 E（ベクトル）。
%%   磁束 Φ のカードは電-7（磁気回路で先に出る）に置かれる見こみなので、この章には置かない（X-2。報告に書いた）。
%%   起電力 \mathcal E のカードは電-4。
%% 出典（11月5日に確かめた）
%%   ・電磁誘導加熱式調理器の利用周波数 20.05 kHz から 100 kHz、高周波出力の定格 10 kW 以下：電波法施行規則 第46条の7 第1項第2号
%%       （e-Gov 法令 API https://laws.e-gov.go.jp/api/1/lawdata/325M50080000014 の本文を curl で読んだ。計画書 D14 と同じ）。
%%   ・誘導式読み書き通信設備（13.56 MHz の誘導電波で記録媒体を読み書き）：同 第44条（同じ本文で確かめた。計画書 D15）。
%%   ・ファラデーが 1831年8月29日に鉄の輪で電磁誘導を見つけた：英国王立研究所（Royal Institution）の所蔵品の説明
%%       （https://rigb.org/node/143 の検索の抜粋。輪の両側にコイルを巻いた）、ペンシルベニア大学図書館の日記の複製の記録
%%       （colenda.library.upenn.edu、Aug. 29th, 1831 の頁。原本は王立研究所）。論文は M. Faraday, "Experimental researches in electricity",
%%       Phil. Trans. R. Soc. 122, 125 (1832)、DOI 10.1098/rstl.1832.0006（計画書 D24 で登録を確かめた）。
%%   ・磁気リコネクション：P. A. Cassak, Y.-H. Liu, M. A. Shay, "A review of the 0.1 reconnection rate problem", J. Plasma Phys. (2017)、
%%       arXiv:1708.03449（本文を curl で読んだ）。電気抵抗だけの考え（Sweet–Parker）では遅すぎることは Parker (1963) から知られ、
%%       観測と計算で速さが約 0.1（規格化した値）になるわけは、まだ理論で説明しきれていない。MMS 衛星（Burch ら 2016）が観測している。
%%       NASA のサイト（science.nasa.gov・mms.gsfc.nasa.gov）は作業環境から届かなかった。
%% 〔要確認〕（確かめられなかった値）
%%   ・銅・アルミニウムの抵抗率 1.68×10⁻⁸・2.65×10⁻⁸ Ω·m（20 ℃）：計画書 C-3 の値（理科年表〔要〕）。電-4 ③ の表の値にそろえる。
%%       NIST の Matula (1979) の再録（srd.nist.gov）は届かなかった。
%%   ・地磁気の上下の成分 3.5×10⁻⁵ T（日本の中ほど）：国土地理院・気象庁地磁気観測所・NOAA の計算サイトは届かなかった。
%%       本文では「とします」と書いた（数学 8-10 の「日本ではおよそ 4.5×10⁻⁵ T」と、伏角が約50°であることから見積もった値）。
%%   ・鉄のなべの導電率 1×10⁷ S/m・比透磁率 100、IC カードのアンテナ 2 μH、リレーのコイル 0.2 H は、例題のための仮の値（本文に「とします」）。
\makeatletter
% 回路の部品（tikzpicture[x=1mm, y=1mm] の中。導線を先にかき、その上に置く。角は導線の向き）
\newcommand\sxiR[2]{\begin{scope}[shift={#1}, rotate=#2]
  \filldraw[fill=white, draw=bkInk, line width=0.7pt] (-4.5,-1.6) rectangle (4.5,1.6);\end{scope}}
\newcommand\sxiC[2]{\begin{scope}[shift={#1}, rotate=#2]
  \fill[white] (-1.0,-3.6) rectangle (1.0,3.6);
  \draw[bkInk, line width=1.1pt] (-1.0,-3.4) -- (-1.0,3.4) (1.0,-3.4) -- (1.0,3.4);\end{scope}}
\newcommand\sxiL[2]{\begin{scope}[shift={#1}, rotate=#2]
  \fill[white] (-5.2,-0.6) rectangle (5.2,3.0);
  \draw[bkInk, line width=0.7pt] (-5,0) arc (180:0:1.25) arc (180:0:1.25) arc (180:0:1.25) arc (180:0:1.25);\end{scope}}
\newcommand\sxiE[2]{\begin{scope}[shift={#1}, rotate=#2]
  \fill[white] (-1.2,-4.2) rectangle (1.2,4.2);
  \draw[bkInk, line width=0.7pt] (-1.0,-4.0) -- (-1.0,4.0);
  \draw[bkInk, line width=1.8pt] (1.0,-2.0) -- (1.0,2.0);\end{scope}}
\newcommand\sxiS[2]{\begin{scope}[shift={#1}, rotate=#2]
  \fill[white] (-3.5,-0.6) rectangle (3.5,3.6);
  \fill[bkInk] (-3,0) circle (0.6) (3,0) circle (0.6);
  \draw[bkInk, line width=0.8pt] (-3,0) -- (2.6,3.0);\end{scope}}
% 検流計（丸に針）  \sxiG{(中心)}{針の角（90 がまん中）}
\newcommand\sxiG[2]{\filldraw[fill=white, draw=bkInk, line width=0.7pt] #1 circle (4);
  \draw[bkRed!80!black, line width=0.9pt] #1 -- ++(#2:3.2); \fill[bkInk] #1 circle (0.5);}
% 棒磁石  \sxibar{(中心)}{角}{長さ}{太さ}  角 0 で右が N
\newcommand\sxibar[4]{\begin{scope}[shift={#1}, rotate=#2]
  \filldraw[fill=bkRed!60, draw=bkInk, line width=0.6pt] (0,{-#4/2}) rectangle ({#3/2},{#4/2});
  \filldraw[fill=bkBlue!45, draw=bkInk, line width=0.6pt] ({-#3/2},{-#4/2}) rectangle (0,{#4/2});
  \node[font=\sffamily\bfseries\footnotesize, text=white] at ({#3/4},0) {N};
  \node[font=\sffamily\bfseries\footnotesize, text=white] at ({-#3/4},0) {S};\end{scope}}
% 横から見たコイル（円筒に巻いた線）  \sxicoil{(左下)}{はば}{高さ}{巻き数}
\newcommand\sxicoil[4]{\begin{scope}[shift={#1}]
  \fill[bkYellow!25] (0,0) rectangle (#2,#3);
  \foreach \sxik in {1,...,#4} {\draw[bkOrange!75!black, line width=0.9pt] ({(\sxik-0.5)*#2/#4},0) -- ({(\sxik-0.5)*#2/#4+0.5*#2/#4},#3);}
  \draw[bkInk, line width=0.5pt] (0,0) rectangle (#2,#3);\end{scope}}
% 紙面の向こうへの磁場（×）と手前への磁場（・）  \sxiinto{(x,y)}  \sxiout{(x,y)}
\newcommand\sxiinto[1]{\draw[bkPurple!80!black, line width=0.6pt] #1 ++(-1.1,-1.1) -- ++(2.2,2.2) #1 ++(-1.1,1.1) -- ++(2.2,-2.2);}
\newcommand\sxiout[1]{\draw[bkPurple!80!black, line width=0.6pt] #1 circle (1.2); \fill[bkPurple!80!black] #1 circle (0.45);}
\tikzset{sxi lab/.style={font=\sffamily\footnotesize, text=bkInk}, sxi dim/.style={bkGray, line width=0.5pt, {Stealth[length=1.8mm]}-{Stealth[length=1.8mm]}},
  sxi arr/.style={line width=1pt, -{Stealth[length=2.4mm]}}, sxi B/.style={bkPurple!80!black, line width=0.7pt, -{Stealth[length=2mm]}}}
\makeatother

\phchapter{8}{電気が生まれる：\\電磁誘導とインダクタンス}{IH調理器は、火がないのになぜなべが熱くなる？}{\bklevelrange{1}{5}}{電気が生まれる：電磁誘導とインダクタンス}
\bkmichishirube{5-12,5-15,8-7,8-8,8-11,8-14,電-4,電-6,電-7}{電-9,電-10,電-12,電-14,電-15}

\begin{goal}
\begin{itemize}[label={\color{bkTheme}●}, leftmargin=1.6\zw, itemsep=0pt, topsep=1pt]
\item \bklevel{1}\enspace 磁石をコイルに出し入れすると電流が流れ、動かす向きと速さで、電流の向きと大きさが変わることを、実験で確かめられる。
\item \bklevel{3}\enspace ファラデーの法則とレンツの法則で、動く棒や発電機に生まれる電圧（誘導起電力）の大きさと向きを求められる。
\item \bklevel{3}\enspace 自己インダクタンス・相互インダクタンスを求め、コイルにたまるエネルギーを計算できる。
\item \bklevel{4}\enspace ファラデーの法則の微分形を読み、磁場のエネルギーから電磁石が鉄を引く力を求められる。
\item \bklevel{5}\enspace 渦電流と表皮効果で、IH調理器や、うすい板を重ねた鉄心のしくみを説明できる。
\end{itemize}
\end{goal}

\phkoushiki{%
\phk{磁石やコイルを動かすと、コイルに電流が流れる（電磁誘導）}{中学理科}{ph:電-8:sec1}
\phk{ファラデーの電磁誘導の法則 $\mathcal E=-N\dfrac{\Delta\Phi}{\Delta t}$、レンツの法則}{高校物理}{ph:電-8:faraday}
\phk{磁場を横切る導体の棒の起電力 $\mathcal E=vBl$}{高校物理}{ph:電-8:vbl}
\phk{交流の発生 $\mathcal E=NBS\omega\sin\omega t$}{高校物理}{ph:電-8:hatsuden}
\phk{自己誘導 $\mathcal E=-L\dfrac{\Delta I}{\Delta t}$、ソレノイド $L=\mu_0\dfrac{N^2S}{l}$}{高校物理}{ph:電-8:L}
\phk{相互誘導 $\mathcal E_2=-M\dfrac{\Delta I_1}{\Delta t}$}{高校物理}{ph:電-8:M}
\phk{コイルのエネルギー $U=\tfrac12LI^2$}{高校物理}{ph:電-8:energy}
}

\begin{hatena}
\noindent\begin{minipage}[c]{0.53\linewidth}\raggedright
IH調理器のなべは、火も見えないし、天板も赤くなっていないのに、お湯がぐらぐらわきます。なべをどけて天板にさわると、なべから熱が移ったところが温かいだけで、天板そのものが熱を出しているわけではありません。

熱は、いったいどこで生まれているのでしょう。そして、鉄のなべはよく熱くなるのに、アルミや銅のなべは熱くなりにくいのは、なぜでしょう。
\end{minipage}\hfill
\begin{minipage}[c]{0.44\linewidth}\centering
\begin{tikzpicture}[x=1mm, y=1mm]
  % なべ
  \fill[bkBlue!20] (8,15) rectangle (48,27);
  \draw[bkInk, line width=1.1pt] (6,32) -- (8,32) -- (8,14) -- (48,14) -- (48,32) -- (50,32);
  \draw[bkInk, line width=0.9pt] (48,29) -- (53,30);
  \foreach \x in {16,24,32,40} {\draw[bkGray, line width=0.5pt] (\x,28.5) .. controls (\x-1.4,31) and (\x+1.4,32.5) .. (\x,35);}
  \fill[bkRed!55] (8,13.2) rectangle (48,14.4);
  % 天板
  \fill[bkGray!30] (0,10.5) rectangle (56,12.6);
  \draw[bkInk, line width=0.5pt] (0,10.5) rectangle (56,12.6);
  % コイル（うずまきを横から）
  \foreach \x in {12,16,...,44} {\draw[bkOrange!80!black, line width=1.1pt] (\x,5.5) ellipse (1.2 and 1.6);}
  \node[sxi lab] at (28,23) {なべ};
  \node[sxi lab, anchor=north] at (28,3) {天板の下のコイル};
  \node[sxi lab, text=bkRed!80!black] at (28,17.3) {底が熱い};
  \node[sxi lab, anchor=south] at (28,36) {火は見えない};
\end{tikzpicture}
\end{minipage}
\end{hatena}

\noindent 第2部の前の2章（電-6・電-7）では、電流が磁場を作ることを見てきました。この章では、その逆向きのこと、つまり\textbf{磁場の変化が電気を生む}ことを調べます。この現象を\textbf{電磁誘導}といいます。\bkidx{でんじゆうどう}{電磁誘導}発電所の発電機も、変圧器も、IH調理器も、スマホのワイヤレス充電も、交通系 IC カードも、電磁誘導で動いています。電気を「つくる」「送る」「使う」の、どの場面にも出てくる、電磁気学の中心の法則の1つです。

\begin{yattemiyou}[磁石とコイルで電流を起こそう]
\textbf{用意するもの}\par
エナメル線（直径\qty{0.2}{mm}くらい、\qty{20}{m}）、トイレットペーパーのしん、紙やすり、テスター（または理科室の検流計）、ネオジム磁石（円柱の形、直径\qty{1}{cm}くらい）、アルミのパイプ（内側の直径が磁石より少し大きく、長さ\qty{1}{m}くらい。ホームセンターにある）、同じくらいの長さの塩化ビニルのパイプ\par\vspace{2pt}
\noindent\begin{minipage}[c]{0.64\linewidth}
\begin{enumerate}[label={\bkcircled[bkGreen]{\arabic*}}, leftmargin=2.4\zw, itemsep=1pt, topsep=0pt]
\item しんにエナメル線を約100回巻き、両はしを紙やすりでみがいて、テスターにつなぐ（テスターは電流の、いちばん小さいはんいにする）。
\item 磁石の N 極を、コイルの穴にすばやく入れる。止める。すばやく出す。テスターの針（数字）を見る。
\item 磁石の向きを逆にして（S 極から）、②をする。ゆっくり動かしたときとも比べる。
\item アルミのパイプを立て、上から磁石を落とす。下に出てくるまでの時間を、ストップウォッチではかる。塩化ビニルのパイプでも同じことをする。
\end{enumerate}
\end{minipage}\hfill
\begin{minipage}[c]{0.33\linewidth}\centering
\begin{tikzpicture}[x=1mm, y=1mm]
  \sxicoil{(4,4)}{22}{12}{11}
  \sxibar{(38,10)}{180}{14}{5}
  \draw[bkInk, sxi arr] (47,16) -- (40,16);
  \draw[bkInk, line width=0.6pt] (6,4) -- (6,-4) -- (12,-4) (24,4) -- (24,-4) -- (20,-4);
  \sxiG{(16,-4)}{70}
  \node[sxi lab, anchor=north] at (16,-8.5) {テスター};
  \node[sxi lab, anchor=south] at (15,16.5) {コイル};
  \node[sxi lab, anchor=south] at (40,17.5) {入れる};
\end{tikzpicture}
\end{minipage}
\par\vspace{2pt}\textbf{気をつけること}：ネオジム磁石はとても強い。2つが急にくっついて指をはさまないように、また、スマホやカードに近づけないようにする。小さい子が口に入れないようにする。
\end{yattemiyou}

\begin{bkcenter}
\begin{bktable}{colspec={Q[l,m,wd=36mm] Q[l,m,wd=22mm] Q[l,m,wd=36mm] Q[l,m,wd=24mm]}, rowsep=1pt, row{1}={bg=bkTheme!10}, cells={font=\small}}
したこと & 針・時間 & したこと & 針・時間 \\
N 極をすばやく入れる & 右へ大きく & S 極をすばやく入れる & 左へ大きく \\
入れたまま止める & 0 にもどる & N 極をゆっくり入れる & 右へ少しだけ \\
N 極をすばやく出す & 左へ大きく & 塩化ビニルのパイプ & 約\qty{0.5}{s}で出る \\
 & & アルミのパイプ & 約\qty{4}{s}かかる \\
\end{bktable}
\bktabcaption{やってみようで、ミナさんが調べた結果}\label{ph:電-8:tab-yattemiyou}
\end{bkcenter}

\noindent やってみようの答え：表\ref{ph:電-8:tab-yattemiyou}から、次のことがわかります。
\begin{enumerate}[label={\bkcircled{\arabic*}}, leftmargin=2.4\zw, itemsep=1pt, topsep=2pt]
\item 磁石が\textbf{動いているあいだだけ}電流が流れ、止めると流れない。電池はどこにもないのに、電流が生まれる。
\item 入れるときと出すときで、電流の向きが反対になる。磁石の極を反対にしても、向きが反対になる。
\item すばやく動かすほど、電流が大きい。
\item アルミは磁石につかないのに、アルミのパイプの中では、磁石が8倍もゆっくり落ちる。塩化ビニルの約\qty{0.5}{s}は、\qty{1}{m}を自由に落ちる時間 $\sqrt{2\times1/9.8}\approx\qty{0.45}{s}$ とほぼ同じです。
\end{enumerate}
①〜③は、この章の③のファラデーの法則で、④は⑪の渦電流で説明します。④のブレーキは、IH調理器でなべが熱くなるのと、もとは同じしくみです。

\begin{omoidasou}
\bkrecall{電-6}{電流の周りには、右ねじの向きに輪の形の磁場ができる。長いソレノイドの中の磁場は $B=\mu_0nI$（$n$ は1 m あたりの巻数）。磁場の単位は T（テスラ）、$\mu_0=4\pi\times10^{-7}$ N/A$^2$。}
\bkrecall{電-6}{電荷 $q$ が速度 $\vec v$ で動くと、磁場から力 $\vec F=q\vec v\times\vec B$（ローレンツ力）を受ける。電流が流れる導線の長さ $l$ の部分は $IlB$ の力を受ける（フレミングの左手）。}
\bkrecall{電-4}{オームの法則 $V=RI$。電流の密度は $\vec j=\sigma\vec E$（$\sigma$ は導電率）。抵抗で出る熱（ジュール熱）は1秒に $RI^2$。}
\bkrecall{電-7}{鉄の中では、同じ電流でも磁場が比透磁率 $\mu_{\mathrm r}$ 倍（数百〜数千倍）になる。透磁率 $\mu=\mu_{\mathrm r}\mu_0$。}
\bkrecall{8-8}{面を通りぬける磁場の量（磁束）は $\Phi=BA\cos\theta$（$\theta$ は磁場と面の向きの角）。発電機は、コイルを通る磁束の変化を使う。}
\bkrecall{8-11}{変わる磁場があると、閉じた道にそって1周する電場の線積分は0にならない：$\oint_C\vec E\cdot\dd\vec r=-\dfrac{\dd\Phi}{\dd t}$（S6、ファラデーの法則）。}
\bkrecall{5-15}{$\dfrac{\dd}{\dd t}\cos\omega t=-\omega\sin\omega t$。}
\end{omoidasou}

\noindent この章の記号：数学の本（8-8）は面積を $A$ と書きましたが、この本の第1〜3部では、高校の教科書と同じく $S$ と書きます。誘導起電力（電磁誘導で生まれる電圧）は、高校の教科書では $V$ と書きますが、この本では電-4 と同じく $\mathcal E$（イー、筆記体）と書きます。この章では、$v$ は棒などの\textbf{速さ}だけに使います。

%% ======================== 物理の道具 ========================
\bksection{doguu}{物理の道具}

\phstage{me}{磁石とコイルで、電気を起こす}{1}

\bksubsection{1}{磁石を動かすと、電流が流れる}
\label{ph:電-8:sec1}
やってみようで見たことを、図\ref{ph:電-8:fig-4koma}にまとめます。コイルの両はしを検流計（小さな電流をはかる計器）につなぎ、棒磁石を動かします。
\begin{bkfigure}
\begin{tikzpicture}[x=1mm, y=1mm]
  \foreach \i/\lab/\ang/\arr/\pole in {0/{N 極を近づける}/45/1/180, 1/{N 極を止める}/90/0/180, 2/{N 極を遠ざける}/135/-1/180, 3/{S 極を近づける}/135/1/0} {
    \begin{scope}[shift={(\i*31.5,0)}]
      \sxicoil{(2,10)}{14}{10}{7}
      \draw[bkInk, line width=0.6pt] (3,10) -- (3,4) -- (5,4) (15,10) -- (15,4) -- (13,4);
      \sxiG{(9,4)}{\ang}
      \sxibar{(23.5,15)}{\pole}{9}{4}
      \ifnum\arr=1 \draw[bkInk, sxi arr] (29,21) -- (23,21); \fi
      \ifnum\arr=-1 \draw[bkInk, sxi arr] (23,21) -- (29,21); \fi
      \node[sxi lab, anchor=north, align=center] at (14,-1.5) {\lab};
    \end{scope}}
\end{tikzpicture}
\bkfigcaption{磁石とコイルと検流計。針は、近づけると右、止めると0、遠ざけると左にふれる。極を反対にすると、近づけても左にふれる。}\label{ph:電-8:fig-4koma}
\end{bkfigure}
大事なのは、「磁石がそこにあること」ではなく、「\textbf{コイルの中の磁場が変わること}」です。磁石を止めてコイルを動かしても、同じように電流が流れます。磁石を使わず、となりのコイルの電流を入れたり切ったりして、磁場を変えても流れます。このように、コイルの中の磁場が変わると、コイルに電圧が生まれ、電流が流れる現象を電磁誘導、そのとき生まれる電圧を\textbf{誘導起電力}、流れる電流を\textbf{誘導電流}といいます。\bkidx{ゆうどうきでんりょく}{誘導起電力}\bkidx{ゆうどうでんりゅう}{誘導電流}

1831年8月29日、イギリスのファラデーは、鉄の輪の両側に2つのコイルを巻き、片方のコイルに電池をつないだり切ったりした瞬間にだけ、もう片方のコイルにつないだ検流計の針がふれることを見つけました（電流が磁場を作ることが知られてから、11年後のことです：電-6 ②）。

\begin{phmijika}[電流が流れるのは「変わっている間」だけ]
電池をつないだままでは、となりのコイルの電流は流れ続けません。ファラデーも、はじめは「強い磁石のそばに置けば、ずっと電流が流れるはず」と考えて実験し、うまくいきませんでした。変化が大事だと気づいたことが、発見のかぎでした。家のコンセントの電気が向きを変え続ける交流なのも、変圧器でこの「変化」を使いたいからです（電-12 ⑤）。
\end{phmijika}

\Needspace{12\baselineskip}\bksubsection{2}{自転車のライトと手回し発電機}
\label{ph:電-8:sec2}
\noindent\begin{minipage}[c]{0.62\linewidth}\raggedright
自転車のライトの発電機（ダイナモ）の中では、タイヤの回転で磁石が回り、その周りに固定したコイルの中の磁場が、N→S→N と変わり続けます。だから、コイルにはずっと電流が流れ、向きは行ったり来たりします（\textbf{交流}）。速くこぐと、磁場の変わり方が速くなり、ライトは明るくなります。
\end{minipage}\hfill
\begin{minipage}[c]{0.35\linewidth}\centering
\begin{tikzpicture}[x=1mm, y=1mm]
  \draw[bkGray, line width=0.6pt] (20,20) circle (17);
  \foreach \a in {0,180} {\begin{scope}[shift={(20,20)}, rotate=\a]
    \fill[bkYellow!30] (13,-5) rectangle (17,5);
    \foreach \y in {-4,-2,...,4} \draw[bkOrange!75!black, line width=0.8pt] (12.5,\y) -- (17.5,\y+0.8);
    \draw[bkInk, line width=0.4pt] (13,-5) rectangle (17,5);\end{scope}}
  \sxibar{(20,20)}{30}{18}{6}
  \draw[bkInk, line width=0.8pt, -{Stealth[length=2mm]}] (20,20) ++(110:10.5) arc (110:160:10.5);
  \node[sxi lab, anchor=north] at (20,1) {回る磁石とコイル};
  \node[sxi lab, anchor=east] at (2.5,28) {コイル};
\end{tikzpicture}
\end{minipage}
\par\vspace{3pt}
理科室の手回し発電機も、中は同じです。豆電球をつないで回すと、何もつながないときより、ハンドルが重くなります。電気のエネルギーは、ただでは生まれません。回す手の仕事が、電気のエネルギーに変わっているのです（④で、このことを式で確かめます）。

\begin{phmijika}[発電所の発電機も、同じしくみ]
火力・水力・原子力の発電所の発電機も、大きな磁石（電磁石）を、蒸気や水の力で回して、周りのコイルに電圧を生ませています。東日本では1秒に50回、西日本では60回、向きが変わる交流です（電-10 ①）。太陽光発電だけは、電磁誘導を使わずに光から直接、電気を作ります（物理シリーズの『原子』で学ぶ）。
\end{phmijika}

\Needspace{12\baselineskip}\phstage{te}{磁束の変わり方から、電圧を求める}{3}

\bksubsection{3}{磁束とファラデーの法則・レンツの法則}
\label{ph:電-8:sec3}
「コイルの中の磁場の変化」を数で表すために、コイルを通りぬける磁場の量を考えます。面積 $S$ の平らな輪を、大きさ $B$ の一様な磁場がつらぬいているとき、
\pheq{\Phi=BS\cos\theta}{Wb ＝ T × m$^2$}
を、輪を通る\textbf{磁束}といいます（8-8 の「磁場の流束」と同じもの）。\bkidx{じそく}{磁束}$\theta$ は、磁場の向きと、輪の面に立てた向き（面の向き $\vec n$）との角です。磁場が面に垂直なら $\Phi=BS$、面にそっていれば $\Phi=0$ です。単位は T·m$^2$ で、これを Wb（ウェーバ）とよびます。磁力線の絵でいえば、磁束は「輪を通りぬける磁力線の本数」にあたります。

ファラデーの実験を詳しく調べると、$N$ 回巻いたコイルに生まれる誘導起電力は、1巻きを通る磁束の、\textbf{変わる速さ}に比例することがわかります。
\pheq{\mathcal E=-N\,\frac{\Delta\Phi}{\Delta t}\qquad\Bigl(\text{変化が速くなる極限では}\ \mathcal E=-N\,\frac{\dd\Phi}{\dd t}\Bigr)}{V ＝ Wb/s}
\label{ph:電-8:faraday}これを\textbf{ファラデーの電磁誘導の法則}といいます。\bkidx{ふぁらでーのでんじゆうどうのほうそく}{ファラデーの電磁誘導の法則}単位は Wb/s＝T·m$^2$/s で、これがちょうど V（ボルト）になります（T＝N/(A·m)、V＝J/C＝N·m/(A·s) から確かめられます）。

\begin{phyomi}[$\mathcal E=-N\,\Delta\Phi/\Delta t$]
「コイルに生まれる電圧は、巻数×（1巻きを通る磁束が、1秒あたりに変わる量）。向きは、変化をじゃまする向き」と読みます。磁束が大きいかどうかは関係がなく、\textbf{変わる速さ}だけで決まります。強い磁石をコイルの中に置いたままでは、電圧は0です。$N$ 倍になるのは、$N$ 個の輪を1列につないだのと同じで、1巻きごとの電圧が足し算されるからです。
\end{phyomi}

\phboxhead{向きを決める：レンツの法則}
式の負の符号は、向きを表します。\textbf{誘導電流は、それが作る磁場が、コイルを通る磁束の変化をじゃまする向きに流れる}。これを\textbf{レンツの法則}といいます。\bkidx{れんつのほうそく}{レンツの法則}
\begin{bkfigure}
\begin{tikzpicture}[x=1mm, y=1mm]
  % 左：近づける
  \begin{scope}
    \draw[bkGray, line width=0.7pt] (20,16) arc (90:-90:3 and 12);
    \draw[bkInk, line width=1.4pt] (20,16) arc (90:270:3 and 12);
    \sxibar{(44,4)}{180}{14}{5}
    \draw[bkInk, sxi arr] (52,12.5) -- (45,12.5);
    \node[sxi lab, anchor=south] at (48,13.5) {近づける};
    \foreach \y in {-2,4,10} \draw[sxi B] (36,\y) -- (25,\y);
    \draw[bkRed!80!black, line width=1.2pt, -{Stealth[length=2.4mm]}] (16.6,9) arc (150:210:10);
    \node[sxi lab, text=bkRed!80!black, anchor=east] at (14.6,4) {誘導電流};
    \draw[bkGreen!60!black, line width=1.2pt, dashed, -{Stealth[length=2.4mm]}] (24,-6) -- (35,-6);
    \node[sxi lab, text=bkGreen!50!black, anchor=north, align=center] at (30,-8) {誘導電流の磁場（破線）\\は磁石を押し返す};
    \node[sxi lab, anchor=south] at (20,17) {コイルの右が N};
  \end{scope}
  % 右：遠ざける
  \begin{scope}[shift={(68,0)}]
    \draw[bkGray, line width=0.7pt] (20,16) arc (90:-90:3 and 12);
    \draw[bkInk, line width=1.4pt] (20,16) arc (90:270:3 and 12);
    \sxibar{(44,4)}{180}{14}{5}
    \draw[bkInk, sxi arr] (45,12.5) -- (52,12.5);
    \node[sxi lab, anchor=south] at (48,13.5) {遠ざける};
    \foreach \y in {-2,4,10} \draw[sxi B] (36,\y) -- (25,\y);
    \draw[bkRed!80!black, line width=1.2pt, -{Stealth[length=2.4mm]}] (16.6,-1) arc (210:150:10);
    \draw[bkGreen!60!black, line width=1.2pt, dashed, -{Stealth[length=2.4mm]}] (35,-6) -- (24,-6);
    \node[sxi lab, text=bkGreen!50!black, anchor=north, align=center] at (30,-8) {誘導電流の磁場（破線）\\は磁石を引きとめる};
    \node[sxi lab, anchor=south] at (20,17) {コイルの右が S};
  \end{scope}
\end{tikzpicture}
\bkfigcaption{レンツの法則（コイルは横から見た輪。濃い線が手前の半分）。左：N 極を近づけると、コイルを左向きに通る磁束（実線）が増えるので、それを打ち消す右向きの磁場（破線）を作る向きに電流が流れる。右：遠ざけると、左向きの磁束が減るので、それを補う左向きの磁場を作る向きに流れる。}\label{ph:電-8:fig-lenz}
\end{bkfigure}
近づけるとコイルは磁石を押し返し、遠ざけると引きとめます。どちらでも、磁石を動かす手は、力にさからって仕事をしなければなりません。その仕事が、電流のエネルギー（抵抗で出る熱）になります。もし向きが逆なら、磁石を少し近づけるだけで、どんどん引きこまれて速くなり、電気も熱もただで生まれてしまいます。レンツの法則は、\textbf{エネルギーの保存}が電磁誘導に課したきまりなのです。

\phboxhead{式の符号の約束}
式で向きを読むときは、次のように約束します。コイルの面の向き $\vec n$ を1つ決め、$\Phi$ はその向きに通る磁束を正とします。起電力 $\mathcal E$ は、$\vec n$ の向きに右ねじを進めるとき、ねじを回す向き（コイルを1周する向き）を正とします（電-6 ② の右ねじと同じ）。$\Phi$ が増えると $\mathcal E<0$、つまり電流は右ねじと反対に回り、$\vec n$ と反対向きの磁場を作ります。これが、レンツの法則の向きです。

\phboxhead{場の言葉で書く}
起電力は、電荷を1周運ぶとき、電場が電荷1 C にする仕事です（8-7）。だから、ファラデーの法則は、閉じた道 $C$ と、それをふちとする面の磁束 $\Phi$ を使って、
\pheq{\oint_C\vec E\cdot\dd\vec r=-\frac{\dd\Phi}{\dd t}}{V ＝ Wb/s}
とも書けます（8-11 の S6）。$N$ 回巻きのコイルは、道が $N$ 回まわるので $N$ 倍になります。止まった電荷が作る電場では、1周の線積分はいつも0でした（電-2 ⑨）。磁場が変わる所では、そうではなくなる、というのがファラデーの法則の中身です（⑨で詳しく読みます）。

\begin{phflow}[ファラデーの法則から、高校の式へ]
\phflowrow{$\displaystyle\oint_C\vec E\cdot\dd\vec r=-\frac{\dd\Phi}{\dd t}$}{道がコイルの線を\\$N$ 回まわる}{$\mathcal E=-N\dfrac{\dd\Phi}{\dd t}$}
\phflowrow{$\mathcal E=-N\dfrac{\dd\Phi}{\dd t}$}{短い時間 $\Delta t$ の\\平均の変化で考える}{$\mathcal E=-N\dfrac{\Delta\Phi}{\Delta t}$}
\end{phflow}

\begin{phbaisuru}{磁石を動かす速さ}
\noindent\begin{minipage}[c]{0.45\linewidth}\raggedright
磁石がコイルを通りぬけるとき、1巻きを通る磁束 $\Phi$（上）と、生まれる起電力 $\mathcal E$（下）の変わり方です。速さを2倍にすると、同じ磁束の変化が半分の時間で起きるので、起電力の山は2倍の高さで、はばは半分になります。山の\textbf{面積}（$\mathcal E$ を時間で積分したもの $=-N\times$（磁束の変化））は変わりません。電圧計の針は大きくふれますが、流れる電気の量（電荷）の合計は、速さによらないのです。
\end{minipage}\hfill
\begin{minipage}[c]{0.53\linewidth}\centering
\begin{tikzpicture}
\begin{axis}[bkplot, name=fa, width=56mm, height=26mm, xmin=-3.2, xmax=3.2, ymin=0, ymax=1.25, xtick=\empty, ytick=\empty, ylabel={$\Phi$（上）}, xlabel={}]
  \addplot[bkcurve, domain=-3.1:3.1, samples=120] {1/(1+(x/0.9)^2)^1.5};
  \addplot[bkcurve2, dashed, domain=-3.1:3.1, samples=120] {1/(1+(x/0.45)^2)^1.5};
\end{axis}
\begin{axis}[bkplot, at={(fa.south)}, anchor=north, yshift=-7mm, width=56mm, height=36mm, xmin=-3.2, xmax=3.2, ymin=-2.6, ymax=2.6, xtick={0}, xticklabels={中央を通る}, ytick={0}, ylabel={$\mathcal E$（下）}, xlabel={$t$}]
  \addplot[bkcurve, domain=-3.1:3.1, samples=160] {3*x/0.81/(1+(x/0.9)^2)^2.5*1.1};
  \addplot[bkcurve2, dashed, domain=-3.1:3.1, samples=160] {3*x/0.2025/(1+(x/0.45)^2)^2.5*1.1};
  \node[font=\sffamily\footnotesize, text=bkBlue!80!black, anchor=west, fill=white, inner sep=1pt] at (axis cs:1.4,0.9) {ふつう（実線）};
  \node[font=\sffamily\footnotesize, text=bkOrange!80!black, anchor=west, fill=white, inner sep=1pt] at (axis cs:0.65,2.05) {2倍（破線）};
\end{axis}
\end{tikzpicture}
\end{minipage}
\end{phbaisuru}

\Needspace{10\baselineskip}\bksubsection{4}{磁場を横切る棒：$vBl$ を2通りに導く}
\label{ph:電-8:sec4}
\noindent\begin{minipage}[c]{0.5\linewidth}\raggedright
2本の平行な金属のレール（間かく $l$）を、片はしで抵抗 $R$ につなぎ、その上に金属の棒をのせます。紙面の向こう向きに、大きさ $B$ の一様な磁場があります。棒を右へ速さ $v$ で動かすと、棒には電圧が生まれ、回路に電流が流れます。この起電力を、2つの見方で求めます。
\end{minipage}\hfill
\begin{minipage}[c]{0.47\linewidth}\centering
\begin{tikzpicture}[x=0.78mm, y=0.78mm]
  \foreach \x in {8,16,24,56,64} \foreach \y in {4,12,20,28} {\sxiinto{(\x,\y)}}
  \draw[bkInk, line width=1pt] (2,0) -- (64,0) (2,32) -- (64,32) (2,0) -- (2,32);
  \sxiR{(2,16)}{90}
  \fill[bkGray!60] (40,-2) rectangle (43,34);
  \draw[bkInk, line width=0.6pt] (40,-2) rectangle (43,34);
  \draw[bkInk, sxi arr] (44,8) -- (54,8);
  \node[sxi lab, anchor=west] at (54.5,8) {$v$};
  \draw[bkRed!80!black, sxi arr] (50,14) -- (50,22);
  \node[sxi lab, anchor=west, text=bkRed!80!black] at (50.5,18) {$I$};
  \draw[bkRed!80!black, sxi arr] (30,35) -- (16,35);
  \draw[bkGreen!50!black, line width=1.2pt, dashed, -{Stealth[length=2.4mm]}] (39,24) -- (29,24);
  \node[sxi lab, anchor=north, text=bkGreen!50!black, fill=white, inner sep=0.8pt] at (33,23.2) {$IlB$};
  \node[sxi lab, anchor=east] at (-2.5,16) {$R$};
  \draw[sxi dim] (69,0) -- (69,32);
  \node[sxi lab, anchor=west] at (69.5,16) {$l$};
  \node[sxi lab, anchor=north] at (33,-2.5) {$\vec B$ は紙面の向こう向き（×）};
\end{tikzpicture}
\end{minipage}

\phboxhead{見方1：棒の中の電荷が受けるローレンツ力}
棒の中の電荷 $q$（正とします）は、棒といっしょに右へ速さ $v$ で動いているので、磁場から $q\vec v\times\vec B$ の力を受けます（電-6 ③）。右向きの $\vec v$ と向こう向きの $\vec B$ の外積は、紙面の上向きです。大きさは $qvB$ で、棒にそって上向きです。電荷を棒のはしからはしまで（長さ $l$）運ぶとき、この力がする仕事は $qvB\times l$ です。電荷1 C あたりの仕事が起電力なので、
\pheq{\mathcal E=vBl}{V ＝ m/s × T × m}
\label{ph:電-8:vbl}です。

\phboxhead{見方2：回路を通る磁束の変化}
棒とレールと抵抗で囲まれた長方形の面積は、1秒に $lv$ ずつ増えます。磁束は、磁場 $B$ と面積のかけ算なので、1秒に $Blv$ ずつ増え、ファラデーの法則から $|\mathcal E|=\dd\Phi/\dd t=vBl$。\textbf{2つの見方で、同じ答え}になりました。向きも、レンツの法則（紙面の向こう向きの磁束が増えるのをじゃまするので、回路の中に手前向きの磁場を作る向き、つまり、図の反時計回り）と、見方1の「棒の中を上向き」とで一致します。

なぜ同じ答えになるのかは、深い問題です。見方1では「磁場の中を電荷が動く」、見方2では「回路の中の磁束が変わる」と、別のことを言っているように見えます。棒といっしょに動く人には棒の中の電荷は止まって見えるので、ローレンツ力ははたらかず、その人には「電場」が見えるはずです。電場と磁場は、見る人の動きによって入れかわる、1つのものの2つの顔なのです（電-14 ⑨。相対性理論の入口）。

\phboxhead{エネルギーで確かめる}
回路の電流は $I=vBl/R$ です。この電流が流れる棒は、磁場から $IlB$ の力を受けます（電-6 ⑧）。向きは、上向きの電流と向こう向きの磁場から左向き、つまり\textbf{動きと反対}です（レンツの法則）。棒を一定の速さで動かし続けるには、手で同じ大きさの力 $IlB$ で押し続けなければなりません。手の仕事率は
\[ IlB\times v=\frac{vBl}{R}\,lB\,v=\frac{(vBl)^2}{R}=RI^2 \]
で、抵抗で出るジュール熱の仕事率 $RI^2$（電-4 ⑤）とぴったり等しくなります。手がした仕事が、電磁誘導を通して、熱に変わったのです。

\begin{phflow}[ローレンツ力から、棒の起電力へ]
\phflowrow{$\vec F=q\vec v\times\vec B$}{棒・速度・磁場が\\たがいに直角}{$\mathcal E=vBl$}
\phflowrow{$\mathcal E=-\dfrac{\dd\Phi}{\dd t}$}{面積が1秒に\\$lv$ ずつ増える}{$\mathcal E=vBl$（同じ）}
\end{phflow}

\begin{phbaisuru}{棒を動かす速さ $v$}
\noindent\begin{minipage}[c]{0.45\linewidth}\raggedright
起電力も電流も2倍になります。棒が受けるブレーキの力 $IlB$ も2倍です。手の仕事率（力×速さ）は $2\times2=4$ 倍、ジュール熱も $RI^2$ で4倍になります。速く動かすほど、ずっと重くなるのです。アルミのパイプの中の磁石が、落ちる速さが増えるにつれてブレーキが強くなり、すぐに一定のゆっくりした速さになったのは、このためです（⑪）。\par\vspace{2pt}{\footnotesize （例題2 の回路：$l=\qty{0.5}{m}$、$B=\qty{0.1}{T}$、$R=\qty{1}{\ohm}$。仕事率は $v^2$ に比例する）}
\end{minipage}\hfill
\begin{minipage}[c]{0.53\linewidth}\centering
\begin{tikzpicture}
\begin{axis}[bkplot, width=48mm, height=40mm, xmin=0, xmax=4.8, ymin=0, ymax=0.075, xtick={1,2,3,4}, ytick={0.01,0.04}, yticklabels={0.01,0.04}, xlabel={$v$（m/s）}, ylabel={仕事率（W）}, scaled y ticks=false]
  \addplot[bkcurve, domain=0:4.5, samples=40] {0.0025*x^2};
  \fill[bkBlue] (axis cs:2,0.01) circle (1.6pt) (axis cs:4,0.04) circle (1.6pt);
  \draw[bkGray, dashed, line width=0.5pt] (axis cs:2,0) -- (axis cs:2,0.01) -- (axis cs:0,0.01) (axis cs:4,0) -- (axis cs:4,0.04) -- (axis cs:0,0.04);
\end{axis}
\end{tikzpicture}
\end{minipage}
\end{phbaisuru}

\Needspace{16\baselineskip}\bksubsection{5}{発電機：回るコイルから交流が生まれる}
\label{ph:電-8:sec5}
\noindent\begin{minipage}[c]{0.46\linewidth}\centering
\begin{tikzpicture}[x=0.86mm, y=0.86mm]
  \fill[bkRed!55] (0,2) rectangle (8,34);
  \fill[bkBlue!40] (52,2) rectangle (60,34);
  \node[font=\sffamily\bfseries\footnotesize, text=white] at (4,18) {N};
  \node[font=\sffamily\bfseries\footnotesize, text=white] at (56,18) {S};
  \foreach \y in {6,12,24,30} \draw[sxi B] (9,\y) -- (51,\y);
  \draw[bkGray, line width=0.5pt, dash pattern=on 2pt off 1pt] (30,-1) -- (30,37);
  \draw[bkOrange!80!black, line width=1.6pt] (22,8) -- (38,12) -- (38,28) -- (22,24) -- cycle;
  \draw[bkInk, line width=1pt, -{Stealth[length=2.2mm]}] (30,18) -- ++(-28:11);
  \node[sxi lab, anchor=west] at (40,15.5) {$\vec n$};
  \draw[bkInk, line width=0.6pt] (38.5,18) arc (0:-28:8.5);
  \node[sxi lab] at (44,20.5) {$\theta$};
  \draw[bkInk, line width=0.8pt, -{Stealth[length=2mm]}] (26,36) arc (160:20:4.5 and 2);
  \node[sxi lab, anchor=south] at (30,38.5) {$\omega$ で回す};
  \node[sxi lab, anchor=north] at (30,0) {面積 $S$、$N$ 巻き};
\end{tikzpicture}
\end{minipage}\hfill
\begin{minipage}[c]{0.5\linewidth}\raggedright
一様な磁場 $B$ の中で、面積 $S$、$N$ 回巻きのコイルを、磁場に垂直な軸の周りに、角速度 $\omega$ で回します。$t=0$ にコイルの面が磁場に垂直（$\vec n$ が磁場の向き）だったとすると、時刻 $t$ の角は $\theta=\omega t$ で、1巻きを通る磁束は
\[\Phi=BS\cos\omega t\]
です。ファラデーの法則で微分すると（5-15）、
\end{minipage}
\pheq{\mathcal E=-N\frac{\dd\Phi}{\dd t}=NBS\omega\sin\omega t}{V ＝ T × m$^2$ × rad/s}
\label{ph:電-8:hatsuden}となります。起電力は正弦波の形で、向きが行ったり来たりする交流です。山の高さ（最大値）は $\mathcal E_{\max}=NBS\omega$、1秒に向きが変わる回数（周波数）は、1秒に回る回数 $f=\omega/(2\pi)$ と同じです。\bkidx{こうりゅう}{交流}

\begin{phyomi}[$\mathcal E=NBS\omega\sin\omega t$]
「磁束がいちばん大きい（コイルの面が磁場に垂直な）とき、起電力は0。磁束が0（面が磁場にそっている）とき、起電力はいちばん大きい」と読みます。磁束の大きさではなく、変わる速さで決まるからです。$\cos$ の山のところでは、$\cos$ のグラフの傾きが0です。磁束のグラフと起電力のグラフは、4分の1周期（90°）ずれています。
\end{phyomi}

\begin{phbaisuru}{回す速さ（角速度 $\omega$）}
\noindent\begin{minipage}[c]{0.45\linewidth}\raggedright
山の高さが2倍になり、同時に、周波数も2倍になります（実線→破線）。自転車のライトが、速くこぐと明るくなるのはこのためです。発電所では、周波数を50 Hz か60 Hz にそろえなければならないので、回す速さを変えずに、電磁石の電流で磁場 $B$ を変えて、電圧を調節しています（電-12 ⑧）。\par\vspace{2pt}{\footnotesize （例題3 の発電機。実線は 50 Hz で山 157 V、破線は2倍の速さで回したとき）}
\end{minipage}\hfill
\begin{minipage}[c]{0.53\linewidth}\centering
\begin{tikzpicture}
\begin{axis}[bkplot, axis lines=left, width=52mm, height=40mm, xmin=0, xmax=21.5, ymin=-360, ymax=360, xtick={5,10,15,20}, ytick={-314,-157,157,314}, xlabel={$t$（ms）}, ylabel={$\mathcal E$（V）}]
  \draw[bkGray, line width=0.5pt] (axis cs:0,0) -- (axis cs:21.5,0);
  \addplot[bkcurve, domain=0:20.5, samples=120] {157.1*sin(18*x)};
  \addplot[bkcurve2, dashed, domain=0:20.5, samples=160] {314.2*sin(36*x)};
\end{axis}
\end{tikzpicture}
\end{minipage}
\end{phbaisuru}

\phboxhead{表計算で、発電機の電圧をかく}
例題3 の発電機（$N=100$、$S=\qty{0.01}{m^2}$、$B=\qty{0.5}{T}$、50 Hz）の起電力を、表計算で $\qty{2.5}{ms}$（45° 回るごと）に計算します。B列は回った角、C列が起電力です。
\begin{bkcenter}
\bksheet{3}{24}{%
1 & $t$（ms） & 角（°） & $\mathcal E$（V） \\
2 & 0 & 0 & 0.0 \\
3 & 2.5 & 45 & 111.1 \\
4 & 5.0 & 90 & 157.1 \\
5 & 7.5 & 135 & 111.1 \\
6 & 10.0 & 180 & 0.0 \\
7 & 12.5 & 225 & $-111.1$ \\
8 & 15.0 & 270 & $-157.1$ \\
9 & 17.5 & 315 & $-111.1$ \\
10 & 20.0 & 360 & 0.0 \\}
\bktabcaption{B列の式は \texttt{=360*50*A2/1000}、C列の式は \texttt{=100*0.5*0.01*2*PI()*50*SIN(RADIANS(B2))}。20 ms で1回まわり、もとにもどる。}\label{ph:電-8:tab-sheet}
\end{bkcenter}
C2 の 50 を 100 に変える（回す速さを2倍にする）と、山が314.2 V になり、10 ms で1回まわるようになります。

\Needspace{14\baselineskip}\bksubsection{6}{自己誘導とインダクタンス}
\label{ph:電-8:sec6}
コイルに流れる電流 $i$ を変えると、コイルが自分で作る磁場が変わるので、コイル自身を通る磁束も変わります。すると、ファラデーの法則で、コイル自身に起電力が生まれます。これを\textbf{自己誘導}といいます。\bkidx{じこゆうどう}{自己誘導}鉄を使わないコイルでは、自分の磁場は電流に比例するので、$N$ 巻きを通る磁束の合計 $N\Phi$ も電流に比例します。
\pheq{N\Phi=Li}{Wb ＝ H × A}
比例定数 $L$ を\textbf{自己インダクタンス}（ふつう、単に\textbf{インダクタンス}）といい、単位は H（ヘンリー）＝Wb/A です。\bkidx{いんだくたんす}{インダクタンス}ファラデーの法則に入れると、
\pheq{\mathcal E=-\frac{\dd(N\Phi)}{\dd t}=-L\frac{\dd i}{\dd t}}{V ＝ H × A/s}
\label{ph:電-8:L}です。電流を増やそうとすると、それをじゃまする向きの起電力が生まれます（レンツの法則）。高校の教科書の $V=-L\,\Delta I/\Delta t$ と同じ式です。

\phboxhead{部品のきまり $V_L=L\,\dd i/\dd t$}
回路にコイルをつないで電流 $i$ を流すとき、コイルの両はしの電圧を「電流の流れこむ側が、流れ出る側より、どれだけ高いか」ではかります（抵抗の電圧 $Ri$ と同じはかり方）。起電力 $-L\,\dd i/\dd t$ は電流をじゃまする向きなので、この電圧は
\[ V_L=L\,\frac{\dd i}{\dd t} \]
です。数学の本（6-4⑦）が「部品のきまり」として使い、「くわしくは物理シリーズ『電磁気学』」と書いた式は、こうしてファラデーの法則から出てくるのです。電流が一定なら $V_L=0$ で、（巻き線の抵抗を無視すれば）コイルはただの導線です。

\phboxhead{長いソレノイドの $L$}
長さ $l$、断面積 $S$、$N$ 回巻きの長いソレノイドに電流 $i$ を流すと、中の磁場は $B=\mu_0(N/l)\,i$ です（電-6 ⑥。1 m あたりの巻数 $n=N/l$）。1巻きを通る磁束は $\Phi=BS$ なので、
\[ N\Phi=N\times\mu_0\frac{N}{l}\,i\times S=\mu_0\frac{N^2S}{l}\,i \]
$N\Phi=Li$ と比べて、
\pheq{L=\mu_0\frac{N^2S}{l}}{H ＝ N/A$^2$ × m$^2$ ÷ m}
です。中に鉄心を入れると、磁場がおよそ $\mu_{\mathrm r}$ 倍（電-7 ③）になるので、$L$ も $\mu_{\mathrm r}$ 倍になります。

\begin{phyomi}[$L=\mu_0N^2S/l$]
「$L$ は、巻数の2乗と断面積に比例し、長さに反比例する。電流によらず、コイルの形だけで決まる」と読みます。巻数が2乗になるのは、巻数を増やすと、①1巻きを通る磁束（磁場）が増え、②その磁束を受けとる巻きの数も増える、の2つが重なるからです。
\end{phyomi}

\begin{phbaisuru}{巻数 $N$}
\noindent\begin{minipage}[c]{0.45\linewidth}\raggedright
長さと太さを変えずに巻数を2倍にすると、$L$ は4倍になります。例題5 のソレノイド（長さ\qty{0.1}{m}・半径\qty{1}{cm}）では、500 回巻きで\qty{0.99}{mH}、1000 回巻きで\qty{3.9}{mH}です。
\end{minipage}\hfill
\begin{minipage}[c]{0.53\linewidth}\centering
\begin{tikzpicture}
\begin{axis}[bkplot, width=52mm, height=40mm, xmin=0, xmax=1250, ymin=0, ymax=6.4, xtick={500,1000}, ytick={0.99,3.95}, yticklabels={0.99,3.9}, xlabel={$N$}, ylabel={$L$（mH）}]
  \addplot[bkcurve, domain=0:1200, samples=50] {3.948e-6*x^2};
  \fill[bkBlue] (axis cs:500,0.987) circle (1.6pt) (axis cs:1000,3.948) circle (1.6pt);
  \draw[bkGray, dashed, line width=0.5pt] (axis cs:500,0) -- (axis cs:500,0.987) -- (axis cs:0,0.987) (axis cs:1000,0) -- (axis cs:1000,3.948) -- (axis cs:0,3.948);
\end{axis}
\end{tikzpicture}
\end{minipage}
\end{phbaisuru}

\begin{phmijika}[ギターのピックアップ、蛍光灯の安定器]
エレキギターの弦の下には、磁石に細い線を数千回巻いたコイル（ピックアップ）があります。鉄の弦がゆれると、コイルを通る磁束が変わり、弦のゆれと同じ形の電圧が生まれます。また、昔の蛍光灯の器具には「安定器」というコイルが入っていて、スイッチを切りかえた瞬間の自己誘導の高い電圧で、蛍光灯に光をつけていました。
\end{phmijika}

\phboxhead{同軸ケーブルの1 m あたりの $L$}
\label{ph:電-8:coax}\noindent\begin{minipage}[c]{0.66\linewidth}\raggedright
テレビのアンテナ線のような同軸ケーブルは、半径 $a$ の中心の線と、それを囲む半径 $b$ の筒でできていて、電流 $I$ が中心の線を行き、筒を帰ってきます。線と筒のあいだ（$a<r<b$）の磁場は、アンペールの法則で、中心の周りの円の道を選ぶと（電-6 ⑥）、
\[ B\times2\pi r=\mu_0I\qquad\text{より}\qquad B=\frac{\mu_0I}{2\pi r} \]
です。筒の外では、行きと帰りの電流が打ち消して、磁場は0です。
\end{minipage}\hfill
\begin{minipage}[c]{0.31\linewidth}\centering
\begin{tikzpicture}[x=1mm, y=1mm]
  \fill[bkGray!35] (16,16) circle (15);
  \fill[white] (16,16) circle (13.2);
  \fill[bkYellow!25] (16,16) circle (13.2);
  \fill[bkOrange!70] (16,16) circle (3.6);
  \draw[bkPurple!80!black, line width=0.7pt, dashed] (16,16) circle (8.5);
  \draw[bkPurple!80!black, line width=0.8pt, -{Stealth[length=2mm]}] (16,16) ++(80:8.5) -- ++(170:0.4);
  \draw[sxi dim] (16,16) -- ++(-40:3.6);
  \node[sxi lab, anchor=north] at (16,0) {断面（$r$ の円の道は破線）};
\end{tikzpicture}
\end{minipage}
\par\vspace{3pt}
長さ $l$ の部分で、線と筒のあいだの長方形（たて $b-a$、横 $l$）を通る磁束は、細い帯（はば $\dd r$）の磁束 $B\,l\,\dd r$ を積分して（5-12）
\[ \Phi=\int_a^b\frac{\mu_0I}{2\pi r}\,l\,\dd r=\frac{\mu_0Il}{2\pi}\ln\frac{b}{a} \]
です。電流と磁束の比 $\Phi/I$ を長さ $l$ でわった、1 m あたりのインダクタンス $L'$ は
\pheq{L'=\frac{\mu_0}{2\pi}\ln\frac{b}{a}}{H/m ＝ N/A$^2$}
となります（中心の線の中の磁場は、周波数が高いと表皮効果（⑫）で電流が表面に集まるので、無視しました）。電-3 ⑩・電-13 の同軸ケーブル（寸法 0.45 mm と 1.6 mm）では、$b/a=1.6/0.45\approx3.56$ で、
\[ L'=2\times10^{-7}\times\ln3.56\approx\qty{2.5e-7}{H/m}=\qty{0.25}{\micro H/m} \]
です。電-13 では、この $L'$ と、1 m あたりの静電容量 $C'$（電-3 ⑩）から、ケーブルの上を進む電気の波の速さと、特性インピーダンスを求めます。

\bksubsection{7}{相互誘導}
\label{ph:電-8:sec7}
\noindent\begin{minipage}[c]{0.58\linewidth}\raggedright
2つのコイルを近くに置き、コイル1の電流 $i_1$ を変えると、コイル1の磁場がコイル2を通る磁束も変わるので、コイル2に起電力が生まれます。これを\textbf{相互誘導}といいます（ファラデーの鉄の輪の実験がこれです）。\bkidx{そうごゆうどう}{相互誘導}コイル2の $N_2$ 巻きを通る磁束の合計は $i_1$ に比例するので、
\[ N_2\Phi_{2}=Mi_1,\qquad \mathcal E_2=-M\frac{\dd i_1}{\dd t} \]
\label{ph:電-8:M}と書けます。$M$ を\textbf{相互インダクタンス}といい、単位は H です。\bkidx{そうごいんだくたんす}{相互インダクタンス}
\end{minipage}\hfill
\begin{minipage}[c]{0.39\linewidth}\centering
\begin{tikzpicture}[x=1mm, y=1mm]
  \draw[bkGray!80, line width=3.5pt] (6,4) rectangle (40,30);
  \foreach \y in {9,12,...,24} \draw[bkOrange!80!black, line width=0.9pt] (2.5,\y) -- (9.5,\y+1.5);
  \foreach \y in {9,11,...,25} \draw[bkGreen!55!black, line width=0.9pt] (36.5,\y) -- (43.5,\y+1);
  \draw[sxi B] (14,30) -- (32,30);
  \draw[sxi B] (32,4) -- (14,4);
  \node[sxi lab, anchor=north] at (23,28.5) {共通の磁束};
  \node[sxi lab, anchor=east] at (1.5,17) {1};
  \node[sxi lab, anchor=west] at (44.5,17) {2};
  \node[sxi lab, anchor=north] at (23,1) {鉄心に巻いた2つのコイル};
\end{tikzpicture}
\end{minipage}

\noindent 同じ長さ $l$・断面積 $S$ の筒に、2つのコイル（$N_1$ 巻きと $N_2$ 巻き）を重ねて巻くと、コイル1が作る磁場 $\mu_0(N_1/l)i_1$ が全部コイル2を通るので、
\pheq{M=\mu_0\frac{N_1N_2S}{l}}{H}
となります。それぞれのコイルの自己インダクタンス $L_1=\mu_0N_1^{\,2}S/l$、$L_2=\mu_0N_2^{\,2}S/l$ と比べると、$M^2=L_1L_2$、つまり $M=\sqrt{L_1L_2}$ です。磁束がもれる（片方のコイルの磁束の一部しか、もう片方を通らない）と、$M$ はこれより小さくなります。

2つのコイルが同じ磁束を共有していれば、1巻きあたりに生まれる電圧は同じなので、
\[ \frac{\mathcal E_2}{\mathcal E_1}=\frac{N_2}{N_1} \]
です。巻数の比で電圧を変えられます。これが\textbf{変圧器}の原理です（電-12 ⑤）。\bkidx{へんあつき}{変圧器}スマホの充電器の中の小さな変圧器も、電柱の上の大きな変圧器も、相互誘導で電圧を上げ下げしています。

\bksubsection{8}{コイルのエネルギーと磁場のエネルギー}
\label{ph:電-8:sec8}
コイルに電流を流しはじめると、自己誘導の起電力が電流をじゃまします。電源は、それにさからって電流を押しこむ仕事をします。電流が $i$ のとき、コイルの電圧は $V_L=L\,\dd i/\dd t$ なので、短い時間 $\dd t$ に電源がする仕事は
\[ V_L\,i\,\dd t=L\frac{\dd i}{\dd t}\,i\,\dd t=Li\,\dd i \]
です。電流を0から $I$ まで増やすあいだの仕事を合計すると（5-12 の積分）、
\pheq{U=\int_0^I Li\,\dd i=\frac12LI^2}{J ＝ H × A$^2$}
\label{ph:電-8:energy}このエネルギーは消えずに、コイルにたまっています。電流を減らすと、コイルはこのエネルギーを返しながら、電流を流し続けようとします。数学の本（6-6）が「部品のきまり」として使った $\tfrac12Li^2$ は、こうして導けます。\bkidx{こいるのえねるぎー}{コイルのエネルギー}

\noindent\begin{minipage}[c]{0.44\linewidth}\centering
\begin{tikzpicture}
\begin{axis}[bkplot, width=58mm, height=42mm, xmin=0, xmax=1.25, ymin=0, ymax=1.3, xtick={1}, xticklabels={$I$}, ytick={1}, yticklabels={$LI$}, xlabel={電流 $i$}, ylabel={$N\Phi$}]
  \fill[bkTheme!25] (axis cs:0,0) -- (axis cs:1,1) -- (axis cs:1,0) -- cycle;
  \addplot[bkcurve, domain=0:1.15] {x};
  \draw[bkGray, dashed, line width=0.5pt] (axis cs:1,0) -- (axis cs:1,1) -- (axis cs:0,1);
  \node[font=\sffamily\footnotesize] at (axis cs:0.68,0.25) {$\tfrac12LI^2$};
\end{axis}
\end{tikzpicture}
\end{minipage}\hfill
\begin{minipage}[c]{0.53\linewidth}\raggedright
$N\Phi$ と電流 $i$ のグラフで見ると、$Li\,\dd i$ は、たての $N\Phi=Li$ と横のはば $\dd i$ の細い長方形です。それを全部たした三角形の面積が $\tfrac12LI^2$ です。コンデンサーの $\tfrac12CV^2$（電-3 ⑦）、ばねの $\tfrac12kx^2$（力-4）と同じ形の三角形です。
\end{minipage}

\phboxhead{エネルギーは、磁場の中にある}
長いソレノイドで、このエネルギーを磁場で書き直してみます。$L=\mu_0N^2S/l$、$B=\mu_0NI/l$（つまり $I=Bl/(\mu_0N)$）を入れると、
\[ U=\frac12\,\mu_0\frac{N^2S}{l}\Bigl(\frac{Bl}{\mu_0N}\Bigr)^2=\frac{B^2}{2\mu_0}\times Sl \]
$Sl$ は、ソレノイドの中の、磁場がある所の体積です。だから、1 m$^3$ あたり
\pheq{u=\frac{B^2}{2\mu_0}}{J/m$^3$ ＝ T$^2$ ÷ (N/A$^2$)}
のエネルギーが、\textbf{磁場のある空間にたまっている}と考えられます。これを\textbf{磁場のエネルギー密度}といいます。\bkidx{じばのえねるぎーみつど}{磁場のエネルギー密度}電場のエネルギー密度 $\tfrac12\varepsilon_0E^2$（電-3 ⑦）と同じ形です。この式は、ソレノイドの中だけでなく、どんな磁場でも成り立ちます（電-14 ⑦）。鉄の中では $B^2/(2\mu)$ で、$\mu$ が大きいので、同じ $B$ でもエネルギーはずっと小さくなります。電磁石のエネルギーは、ほとんどがすき間の空気の中にたまっています（電-7 ⑥）。

\begin{phbaisuru}{電流 $I$}
エネルギーは $I^2$ に比例するので、4倍になります。磁場 $B$ も2倍で、$B^2$ に比例するエネルギー密度は4倍。2つの見方は、きちんと合っています。
\end{phbaisuru}

\begin{motto}[4]
\textbf{$M$ は、どちらから数えても同じ}\par
2つのコイルに電流 $I_1$・$I_2$ を流しておくと、たまるエネルギーは
\[ U=\tfrac12L_1I_1^{\,2}+\tfrac12L_2I_2^{\,2}+MI_1I_2 \]
です。はじめにコイル1を $I_1$ まで、次にコイル2を $I_2$ まで増やすと、2つめのとき、コイル1にはコイル2から $M_{12}\,\dd i_2/\dd t$ の電圧がかかり、電源1は $M_{12}I_1I_2$ だけ余分に仕事をします。逆の順に増やすと、余分の仕事は $M_{21}I_1I_2$ です。たまったエネルギーが増やす順によって違えば、片方の順でためて別の順で取り出すだけで、エネルギーが生まれてしまいます。だから $M_{12}=M_{21}$ で、「1から2へ」の $M$ と「2から1へ」の $M$ は、コイルの形がどんなに違っても等しいのです。
\end{motto}

\begin{phmodel}[コイルとソレノイド]
\phignore{コイルのはしの近くで磁場が弱くなること（はしの効果）。巻き線の抵抗と太さ。コイルの外の磁場。}
\phassume{コイルは十分に長く、中の磁場は一様。鉄心があるときは $B=\mu H$ の直線（電-7 ③）で、磁束はもれない。}
\phlimit{短いコイル（$L$ は式より小さくなる）。鉄心が飽和するほどの電流（$L$ が急に小さくなる：電-7 ⑧）。周波数が高く、コイルの線どうしのあいだの静電容量が効くとき。}
\end{phmodel}

\phkokomade

\Needspace{9\baselineskip}\phstage{shiki}{うずの電場と渦電流を読む}{5}

\bksubsection{9}{微分形：変わる磁場は、電場のうずを作る}
\label{ph:電-8:sec9}
③の積分形 $\oint_C\vec E\cdot\dd\vec r=-\dd\Phi/\dd t$ は、「道 $C$ を選んで計算する」ための形でした。これを、場所ごとの式に直します。ストークスの定理（8-11）で左辺を、道をふちとする面の上の積分にし、右辺の磁束も同じ面の積分 $\Phi=\iint\vec B\cdot\vec n\,\dd S$ で書くと、
\[ \iint_S(\mathrm{rot}\,\vec E)\cdot\vec n\,\dd S=-\iint_S\frac{\partial\vec B}{\partial t}\cdot\vec n\,\dd S \]
これが、どんなに小さな面でも成り立つので、
\pheq{\mathrm{rot}\,\vec E=-\frac{\partial\vec B}{\partial t}}{V/m$^2$ ＝ T/s}
\label{ph:電-8:rot}です（8-14 の4つの式の（3））。$\mathrm{rot}$ は「その点の周りを小さく1周したときの、うずの強さ」でした（8-10）。止まった電荷が作る電場は $\mathrm{rot}\,\vec E=\vec0$（電-2 ⑨）で、うずがありませんでした。\textbf{時間とともに変わる磁場は、その周りに電場のうずを作る}のです。この電場は、電荷がなくても、導線がなくても、空間にできます。

\noindent\begin{minipage}[c]{0.42\linewidth}\centering
\begin{tikzpicture}[x=1mm, y=1mm]
  \fill[bkPurple!10] (22,22) circle (11);
  \draw[bkPurple!60!black, line width=0.6pt, dashed] (22,22) circle (11);
  \foreach \p in {(17,22),(27,22),(22,17),(22,27),(22,22)} {\sxiout{\p}}
  \foreach \r in {6,14,19} {\draw[bkRed!75!black, line width=0.9pt] (22,22) circle (\r);
    \draw[bkRed!75!black, line width=0.9pt, -{Stealth[length=2.2mm]}] (22,22) ++(95:\r) -- ++(5:0.3);}
  \node[sxi lab, anchor=north west] at (36,6) {$\vec E$（実線）};
  \node[sxi lab, anchor=south, align=center] at (22,42) {手前向きの $B$ が増える};
  \draw[sxi dim] (22,22) -- ++(-30:11);
  \node[sxi lab, anchor=north east] at (29,18) {$a$};
\end{tikzpicture}
\end{minipage}\hfill
\begin{minipage}[c]{0.55\linewidth}\raggedright
半径 $a$ の円の中だけに一様な磁場 $B$ があり、時間とともに増えるとき、円と同じ中心の、半径 $r$ の円の道を選びます（8-11 の S6）。対称なので、電場は道にそって大きさが一定です。$r<a$ では
\[ E\times2\pi r=\pi r^2\frac{\dd B}{\dd t},\qquad E=\frac{r}{2}\,\frac{\dd B}{\dd t} \]
$r>a$ では磁束は $\pi a^2B$ のままなので、$E=\dfrac{a^2}{2r}\,\dfrac{\dd B}{\dd t}$ です。磁場のない所にも、電場のうずはとどきます。向きは、レンツの法則のとおり、増える磁場を打ち消す向きの電流を流そうとする向き（図では時計回り）です。
\end{minipage}

\phboxhead{電圧が「道しだい」になる}
電場にうずがあると、2点のあいだの電圧（電場の線積分）が、道によって違ってきます。電位 $\phi$ を1つに決めることはできません。それでも、ふつうの回路の計算でキルヒホッフの電圧則（電-4 ⑦）が使えるのは、変わる磁場がコイルの中にほぼ閉じこめられていて、コイルの外の導線や部品の周りでは $\mathrm{rot}\,\vec E\approx\vec0$ だからです。コイルの中の変わる磁束の効果は、コイルの両はしの電圧 $V_L=L\,\dd i/\dd t$ としてまとめて数えます。電-9 ②では、この約束で回路の式を立てます。

\begin{phmijika}[電子を加速する：ベータトロン]
ドーナツの形の真空の管の中で、電子を円の上に走らせ、管の中心を通る磁場を増やしていくと、電子は電場のうずに押され続けて、どんどん速くなります。この装置をベータトロンといい、医療や研究で電子を加速するのに使われてきました。導線がなくても、変わる磁場は電荷を動かすことができるのです。
\end{phmijika}

\Needspace{12\baselineskip}\bksubsection{10}{電磁石が鉄を引く力}
\label{ph:電-8:sec10}
\noindent\begin{minipage}[c]{0.4\linewidth}\centering
\begin{tikzpicture}[x=1mm, y=1mm]
  \fill[bkGray!45] (4,16) -- (4,34) -- (40,34) -- (40,16) -- (32,16) -- (32,26) -- (12,26) -- (12,16) -- cycle;
  \draw[bkInk, line width=0.6pt] (4,16) -- (4,34) -- (40,34) -- (40,16) -- (32,16) -- (32,26) -- (12,26) -- (12,16) -- cycle;
  \foreach \y in {27,29,...,33} {\draw[bkOrange!80!black, line width=0.9pt] (13,\y) -- (31,\y+1);}
  \fill[bkGray!70] (2,4) rectangle (42,10);
  \draw[bkInk, line width=0.6pt] (2,4) rectangle (42,10);
  \draw[sxi dim] (46,10) -- (46,16);
  \node[sxi lab, anchor=west] at (46.5,13) {$x$};
  \foreach \xx in {6,8,10,34,36,38} \draw[sxi B] (\xx,15.6) -- (\xx,10.4);
  \draw[bkInk, sxi arr] (22,6) -- (22,14.5);
  \node[sxi lab, anchor=west] at (23,13.2) {$F$};
  \node[sxi lab, anchor=north] at (22,3) {鉄の板};
  \node[sxi lab, anchor=south] at (22,35) {電磁石（すき間2つ）};
\end{tikzpicture}
\end{minipage}\hfill
\begin{minipage}[c]{0.57\linewidth}\raggedright
電磁石の極と鉄の板のあいだに、面積 $S$、はば $x$ のすき間があり、すき間の磁場が $B$ だとします。すき間1つにたまっている磁場のエネルギーは、$\dfrac{B^2}{2\mu_0}\times Sx$ です（⑧。鉄の中のエネルギーは、$\mu$ が大きいので無視します）。

電磁石のつくる磁束を一定に保ったまま、力 $F$ にさからって、板を $\dd x$ だけゆっくり引きはなすと、すき間の体積が $S\,\dd x$ 増え、磁場のエネルギーが増えます。
\end{minipage}
\[ F\,\dd x=\frac{B^2}{2\mu_0}\,S\,\dd x\qquad\text{より}\qquad F=\frac{B^2S}{2\mu_0} \]
\label{ph:電-8:kyuuin}引きはなすのにした仕事が、増えたエネルギーになる、と考えたのです。このように、物を少しだけ動かしたと考えて、エネルギーの変化から力を求める方法を\textbf{仮想変位}（仮想仕事）の方法といいます（電-3 ⑪ のコンデンサーの極板の引力と同じ方法）。\bkidx{かそうへんい}{仮想変位}

\begin{phyomi}[$F/S=B^2/(2\mu_0)$]
「磁力線は、1 m$^2$ あたり $B^2/(2\mu_0)$ の力で、ゴムひものように縮もうとする」と読めます。$B=\qty{1}{T}$ なら $1/(2\times4\pi\times10^{-7})\approx\qty{4.0e5}{N/m^2}$ で、大気の圧力の約4倍です。磁場が2倍になれば、引く力は4倍です。
\end{phyomi}

\noindent 電流を一定に保って引きはなすと考えても、同じ力になります。ただし、そのときは電源も仕事をするので、エネルギーの数え方が変わります（専門の本で扱います）。

\Needspace{12\baselineskip}\bksubsection{11}{渦電流：金属の板の中を回る電流}
\label{ph:電-8:sec11}
金属の板の近くで磁場が変わると、⑨の電場のうずが板の中にもでき、板の中を、うずをまく電流が流れます。これを\textbf{渦電流}といいます。\bkidx{うずでんりゅう}{渦電流}渦電流は、板の抵抗でジュール熱を出し（\textbf{渦電流損}）、レンツの法則で磁場の変化をじゃまします。\bkidx{うずでんりゅうそん}{渦電流損}
\begin{itemize}[leftmargin=1.2\zw, itemsep=1pt, topsep=2pt, label={\color{bkTheme}●}]
\item アルミのパイプの中を磁石が落ちると、磁石の上と下で、パイプを通る磁束が変わり、パイプの輪に渦電流が流れます。その磁場が、磁石の落下にブレーキをかけます（やってみよう④）。大型トラックやバスの補助ブレーキにも、回る金属の円板に渦電流を流して止める装置があります。
\item 変圧器やモーターの鉄心では、渦電流は電気のむだ（熱）になります。そこで、鉄心を、表面を絶縁したうすい鉄の板（電磁鋼板）を重ねて作ります。
\end{itemize}

\begin{bkfigure}
\begin{tikzpicture}[x=1mm, y=1mm]
  \fill[bkGray!35] (0,0) rectangle (40,26);
  \draw[bkInk, line width=0.6pt] (0,0) rectangle (40,26);
  \foreach \s in {3,7,10.5} \draw[bkRed!75!black, line width=0.9pt] ({\s},{\s*0.6}) rectangle ({40-\s},{26-\s*0.6});
  \draw[bkRed!75!black, line width=0.9pt, -{Stealth[length=2mm]}] (20,1.8) -- (21,1.8);
  \node[sxi lab, anchor=north] at (20,-2) {1枚のかたまり：大きな渦電流};
  \begin{scope}[shift={(66,0)}]
    \foreach \i in {0,...,7} {\fill[bkGray!35] ({\i*5},0) rectangle ({\i*5+4.4},26);
      \draw[bkInk, line width=0.5pt] ({\i*5},0) rectangle ({\i*5+4.4},26);
      \draw[bkRed!75!black, line width=0.8pt] ({\i*5+1.2},2) rectangle ({\i*5+3.2},24);}
    \node[sxi lab, anchor=north] at (20,-2) {うすい板を重ねる：渦電流は小さい};
  \end{scope}
  \foreach \p in {(46,3),(46,23),(56,3),(56,23)} {\sxiout{\p}}
  \node[sxi lab, align=center] at (51,13) {変わる\\磁場};
\end{tikzpicture}
\bkfigcaption{鉄心の断面。磁場は紙面に垂直で、時間とともに変わる。うずの電場が回る道の長さは、板が厚いほど長く、囲む磁束も大きいので、渦電流損は板の厚さの2乗に比例する（実線が渦電流の道）。}\label{ph:電-8:fig-sekisou}
\end{bkfigure}

\phboxhead{うすい板の渦電流損を式で求める}
厚さ $d$ の広い板の中に、板の面にそった一様な磁場 $B=B_{\mathrm m}\sin\omega t$ があるとします。板のまん中から $x$ だけ離れた所の電場は、⑨と同じ考えで $E=x\,\dd B/\dd t$ です（まん中の面をはさんで、反対向きに流れる）。電流の密度は $j=\sigma E$（電-4 ③）、1 m$^3$ あたりのジュール熱は $jE=\sigma E^2$ です。板の厚さで平均し、時間で平均すると（$\cos^2$ の平均は $\tfrac12$：4-9）、
\[ \frac1d\int_{-d/2}^{d/2}\sigma x^2\,\dd x\times\overline{\Bigl(\frac{\dd B}{\dd t}\Bigr)^2}=\frac{\sigma d^2}{12}\times\frac{\omega^2B_{\mathrm m}^{\,2}}{2} \]
$\omega=2\pi f$ を入れて、
\pheq{P_{\text{渦}}=\frac{\pi^2\sigma f^2B_{\mathrm m}^{\,2}d^2}{6}}{W/m$^3$}
\label{ph:電-8:uzuson}となります（1 m$^3$ あたりの平均の仕事率）。ただし、渦電流が自分で作る磁場が、もとの磁場を変えないくらい、板がうすいときの式です（板の厚さが、⑫の表皮深さより十分小さいとき）。

\begin{phyomi}[$P_{\text{渦}}\propto\sigma f^2B_{\mathrm m}^{\,2}d^2$]
「渦電流損は、周波数の2乗、磁場の2乗、板の厚さの2乗に比例する」と読みます。板の厚さを半分にすると $\tfrac14$ です。周波数の2乗に比例するので、周波数が高いモーターや変圧器ほど、うすい板や、電気を通しにくい材料（フェライト）を使います。導電率 $\sigma$ に比例するので、鉄にけい素をまぜて電気を通しにくくした「けい素鋼板」が使われます。
\end{phyomi}

\phboxhead{プログラムで、板の厚さを変えてみる}
電磁鋼板の導電率を $\sigma=\qty{2e6}{S/m}$、密度を $\qty{7650}{kg/m^3}$ とし、50 Hz・山 1.5 T で、板の厚さを変えます。
\begin{phcode}[Python]
import math
sig, f, Bm, rho = 2.0e6, 50, 1.5, 7650   \# 導電率、周波数、磁場の山、密度
for dmm in [1.0, 0.5, 0.35, 0.23]:      \# 板の厚さ（mm）
    d = dmm / 1000
    p = math.pi**2 * sig * f**2 * Bm**2 * d**2 / 6   \# W/m3
    print(dmm, round(p), round(p / rho, 2))       \# 1 kg あたりは W/kg
\end{phcode}
\begin{phcodeout}
1.0 18506 2.42
0.5 4626 0.6
0.35 2267 0.3
0.23 979 0.13
\end{phcodeout}
左から、板の厚さ（mm）、1 m$^3$ あたり（W）、1 kg あたり（W）です。\qty{1}{mm}の板を\qty{0.35}{mm}にすると、渦電流損は約8分の1になります。実際の鉄心には、このほかにヒステリシス損（電-7 ⑧）もあります。

\Needspace{14\baselineskip}\bksubsection{12}{表皮効果と IH調理器}
\label{ph:電-8:sec12}
周波数が高くなると、金属の中の渦電流は、もとの磁場を打ち消す磁場を強く作るので、磁場も電流も、金属の\textbf{表面近くにしか入れなく}なります。これを\textbf{表皮効果}といいます。\bkidx{ひょうひこうか}{表皮効果}

\phboxhead{磁場は、金属の中へ「しみこむ」}
金属の中では、電場は $\vec j=\sigma\vec E$ の電流を流します。⑨の $\mathrm{rot}\,\vec E=-\partial\vec B/\partial t$ と、アンペールの法則 $\mathrm{rot}\,\vec B=\mu\vec j$（電-6 ⑦。金属の中では透磁率 $\mu$）を組み合わせると（8-12 の $\mathrm{rot}\,\mathrm{rot}=\mathrm{grad}\,\mathrm{div}-\nabla^2$ と $\mathrm{div}\,\vec B=0$ を使う）、
\[ \frac{\partial\vec B}{\partial t}=\frac{1}{\mu\sigma}\,\nabla^2\vec B \]
が出ます。これは、熱が物の中にしみこむ熱伝導の式や、拡散の式（8-14⑧）と同じ形です。磁場は、金属の中へ、熱のように「しみこんで」いきます。平らな金属の表面（深さ $z=0$）で、磁場を $B_0\cos\omega t$ でゆらすと、答えは
\[ B=B_0\,e^{-z/\delta}\cos\Bigl(\omega t-\frac{z}{\delta}\Bigr),\qquad \delta=\sqrt{\frac{2}{\omega\mu\sigma}} \]
です。$B=\mathrm{Re}\bigl[B_0e^{j\omega t-(1+j)z/\delta}\bigr]$（6-10）と書いて式に入れると、左辺は $j\omega$ 倍、右辺は $(1+j)^2/(\mu\sigma\delta^2)=2j/(\mu\sigma\delta^2)$ 倍で、$\delta^2=2/(\omega\mu\sigma)$ なら等しくなることが確かめられます。
\pheq{\delta=\sqrt{\frac{2}{\omega\mu\sigma}}}{m ＝ $\sqrt{1\div(\text{1/s}\times\text{H/m}\times\text{S/m})}$}
\label{ph:電-8:delta}$\delta$ を\textbf{表皮深さ}といいます。\bkidx{ひょうひふかさ}{表皮深さ}深さ $\delta$ ごとに、ゆれの大きさが $e^{-1}\approx0.37$ 倍に減ります。電流も表面から $\delta$ くらいの層だけを流れます。

\begin{phbaisuru}{周波数 $f$}
\noindent\begin{minipage}[c]{0.45\linewidth}\raggedright
$\delta$ は $1/\sqrt2\approx0.71$ 倍になります。周波数が100倍なら、$\delta$ は10分の1です。銅では、コンセントの50 Hz で約\qty{9}{mm}（電-15 ⑩）なので、ふつうの電線の中は、ほぼ一様に電流が流れます。IH の25 kHz では約\qty{0.4}{mm}です。\par\vspace{2pt}{\footnotesize （上から、アルミ（点線）・銅（実線）・鉄（破線）。鉄は導電率 $\qty{1e7}{S/m}$、比透磁率100 とした。どの線も $\sqrt f$ に反比例する）}
\end{minipage}\hfill
\begin{minipage}[c]{0.53\linewidth}\centering
\begin{tikzpicture}
\begin{axis}[bkplot, axis lines=left, xmode=log, ymode=log, width=50mm, height=46mm, xmin=30, xmax=200000, ymin=0.03, ymax=20, xtick={100,1000,10000,100000}, xticklabels={100,1000,$10^4$,$10^5$}, ytick={0.1,1,10}, yticklabels={0.1,1,10}, xlabel={$f$（Hz）}, ylabel={$\delta$（mm）}]
  \addplot[bkcurve, domain=40:150000, samples=40] {9.2255*sqrt(50/x)};
  \addplot[bkcurve3, densely dotted, domain=40:150000, samples=40] {11.5867*sqrt(50/x)};
  \addplot[bkcurve2, dashed, domain=40:150000, samples=40] {2.2508*sqrt(50/x)};
  \draw[bkGray, line width=0.5pt] (axis cs:25000,0.03) -- (axis cs:25000,20);
  \node[font=\sffamily\footnotesize, anchor=east, fill=white, inner sep=1pt] at (axis cs:22000,8) {IH};
\end{axis}
\end{tikzpicture}
\end{minipage}
\end{phbaisuru}

\phboxhead{\bkhatenaref の答え}
\label{ph:電-8:hatena}IH調理器の天板の下には、うずまきの形のコイルがあり、1秒に2万〜10万回ほど向きが変わる電流を流しています（家庭の IH は、電波法施行規則で、20.05 kHz から100 kHz のあいだの周波数を使うと決められています）。コイルの磁場は天板を通りぬけ、なべの底を通る磁束が、激しく変わり続けます。すると、なべの底に渦電流が流れ（⑪）、なべの金属の抵抗でジュール熱が出ます。\textbf{熱は、なべの底そのものの中で生まれている}のです。天板はガラスなどの電気を通さない物なので、渦電流は流れず、自分では熱くなりません。

鉄のなべがよく熱くなるのは、表皮効果のためです。25 kHz の表皮深さは、鉄（比透磁率を100 とします）では約\qty{0.10}{mm}、銅では約\qty{0.41}{mm}、アルミでは約\qty{0.52}{mm}です（例題10）。電流が流れる層の厚さを $\delta$ とみると、その層の電気の流れにくさは $1/(\sigma\delta)$ に比例します。鉄は、導電率 $\sigma$ が銅の約6分の1で、しかも層がうすいので、流れにくさは銅の約24倍です。同じ大きさの渦電流なら、ジュール熱 $RI^2$ も約24倍になります。銅やアルミは電気を通しやすすぎて、渦電流は流れても、熱があまり出ないのです。鉄では、さらにヒステリシス損（電-7 ⑧）も熱になります。このため、IH の多くは、鉄や、磁石につくステンレスのなべを使うように作られています。

\Needspace{14\baselineskip}\bksubsection{13}{IC カードとワイヤレス充電：相互誘導と共振}
\label{ph:電-8:sec13}
\noindent\begin{minipage}[c]{0.56\linewidth}\raggedright
交通系 IC カードには、電池が入っていません。カードのふちにそって、平らなコイル（アンテナ）が数回巻いてあり、小さなコンデンサーと IC チップがつながっています。読みとり機のコイルには、13.56 MHz（1秒に1356万回）で向きが変わる電流が流れていて、その磁場がカードのコイルを通ると、相互誘導（⑦）で起電力が生まれます。これがチップを動かす電気になります。
\end{minipage}\hfill
\begin{minipage}[c]{0.41\linewidth}\centering
\begin{tikzpicture}[x=1mm, y=1mm]
  \fill[bkGray!25] (0,0) rectangle (44,9);
  \draw[bkInk, line width=0.5pt] (0,0) rectangle (44,9);
  \foreach \s in {1.5,3,4.5} \draw[bkOrange!80!black, line width=0.8pt] (6+\s,6) -- (38-\s,6);
  \node[sxi lab] at (22,2.8) {読みとり機のコイル};
  \fill[bkBlue!12, rounded corners=1mm] (4,18) rectangle (40,30);
  \draw[bkInk, line width=0.5pt, rounded corners=1mm] (4,18) rectangle (40,30);
  \foreach \s in {1,2,3} \draw[bkGreen!55!black, line width=0.7pt, rounded corners=0.6mm] (4+\s,18+\s) rectangle (40-\s,30-\s);
  \fill[bkInk!70] (19,22.5) rectangle (25,25.5);
  \node[sxi lab, anchor=south] at (22,30.5) {カード（コイルと IC）};
  \foreach \x in {12,22,32} \draw[sxi B, <->] (\x,10) -- (\x,17);
\end{tikzpicture}
\end{minipage}
\par\vspace{3pt}
起電力を大きくするために、カードのコイルとコンデンサーは、13.56 MHz で\textbf{共振}するように選んであります。コイルとコンデンサーの回路は、決まった周波数 $f=1/(2\pi\sqrt{LC})$ でゆれやすい性質があり（電気振動：電-9 ⑦）、その周波数の磁場を受けると、大きな電流がゆれるのです（共振：電-10 ⑦）。13.56 MHz は、電波法施行規則で「誘導式読み書き通信設備」に決められた周波数です。スマホのワイヤレス充電も、充電台のコイルとスマホの中のコイルの相互誘導で電気を送ります。

%% ---------------------------------------------------------------
\begin{phmada}[磁力線は、なぜあんなに速くつなぎかわる？]
太陽の表面の爆発（フレア）や、オーロラのもとになる地球の周りの宇宙では、向きの反対な磁力線が近づくと、とつぜん磁力線がつなぎかわり、磁場のエネルギーが一気に、熱と、粒子の速い運動に変わります。これを\textbf{磁気リコネクション}といいます。宇宙のガス（プラズマ）は電気をとてもよく通すので、ファラデーの法則と電気抵抗だけで計算すると、つなぎかわりには、実際の何けたも長い時間がかかってしまいます。このことは1960年代から知られていました。計算機のシミュレーションや、地球の周りを回る人工衛星（アメリカの MMS など）の観測で、つなぎかわる場所の細かいようすがわかってきましたが、速さがなぜ観測される値になるのかは、まだ理論で説明しきれていません。
\end{phmada}

\begin{phsaki}{電磁誘導の先へ}
電磁誘導は、モーター・発電機・変圧器などの電気機器（電-12）の土台です。電気機器の設計では、鉄心の中の渦電流や、複雑な形のコイルの磁場を、コンピュータで細かく区切って計算します（有限要素法による電磁界解析）。ファラデーの法則とアンペールの法則に「変位電流」を加えると、電場と磁場が互いに生み合って空間を進む電磁波が出てきます（電-14・電-15）。抵抗が0になる超伝導のコイルでは、電流が流れ続け、強い磁場を作れます（しくみは、物理シリーズの『原子』で学びます）。
\end{phsaki}

\begin{senmonsho}[電気機器・電気回路の専門書の書き方]
\begin{bkhonbun}
$N$ 巻のコイルを貫く磁束を $\phi$ とすれば，コイルと鎖交する磁束は $\psi=N\phi$ であり，これを鎖交磁束（磁束鎖交数）という．電磁誘導によりコイルに誘導される起電力は
\[ e=-\frac{d\psi}{dt}=-N\frac{d\phi}{dt}\ \text{[V]} \]
である．自己インダクタンスは $L=\psi/i$ [H] で定義され，$e=-L\,di/dt$ を逆起電力という．
\end{bkhonbun}
\bkyaku{「鎖交」は、くさりの輪のようにからみ合うこと。$\psi$（プサイ）は、この章の $N\Phi$。専門書は、時間で変わる起電力を小文字の $e$、磁束を小文字の $\phi$ で書くことが多い（この本は $\mathcal E$・$\Phi$。この本の $\phi$ は電位）。「逆起電力」は、電流の変化に逆らう起電力のこと。電気工学の本は、電場を「電界」、磁場（$\vec H$）を「磁界」とよぶ。}
\end{senmonsho}

\begin{taisetsu}
\def\tsl#1{\makebox[7.2\zw][l]{\sffamily\bfseries\small\color{\bkTT}#1}}
\tsl{磁束}$\Phi=BS\cos\theta$〔Wb〕\\[3pt]
\tsl{ファラデー}$\mathcal E=-N\dfrac{\dd\Phi}{\dd t}$、$\displaystyle\oint_C\vec E\cdot\dd\vec r=-\frac{\dd\Phi}{\dd t}$、$\mathrm{rot}\,\vec E=-\dfrac{\partial\vec B}{\partial t}$\\[3pt]
\tsl{動く棒・発電機}$\mathcal E=vBl$、$\mathcal E=NBS\omega\sin\omega t$\\[3pt]
\tsl{インダクタンス}$N\Phi=Li$、$V_L=L\dfrac{\dd i}{\dd t}$、$L=\mu_0\dfrac{N^2S}{l}$、$\mathcal E_2=-M\dfrac{\dd i_1}{\dd t}$\\[3pt]
\tsl{エネルギー}$U=\tfrac12LI^2$、$u=\dfrac{B^2}{2\mu_0}$、吸引力 $F=\dfrac{B^2S}{2\mu_0}$\\[3pt]
\tsl{渦電流・表皮}$P_{\text{渦}}=\dfrac{\pi^2\sigma f^2B_{\mathrm m}^{\,2}d^2}{6}$、$\delta=\sqrt{\dfrac{2}{\omega\mu\sigma}}$
\end{taisetsu}

\begin{kotoba}
\small
\bkword{電磁誘導}{コイルを通る磁束が変わると、コイルに電圧（誘導起電力）が生まれること。}
\bkword{磁束}{面を通りぬける磁場の量。単位 Wb（ウェーバ）＝T·m$^2$。}
\bkword{レンツの法則}{誘導電流は、磁束の変化をじゃまする向きに流れる。}
\bkword{自己インダクタンス}{コイルの電流と、それが自分を通る磁束（の巻数倍）との比 $L$。単位 H（ヘンリー）。}
\bkword{相互インダクタンス}{一方のコイルの電流と、それがもう一方のコイルを通る磁束（の巻数倍）との比 $M$。}
\bkword{渦電流}{変わる磁場のそばの金属の中を、うずをまいて流れる電流。}
\bkword{表皮効果}{周波数が高いと、電流が金属の表面近くだけを流れること。その深さが表皮深さ $\delta$。}
\end{kotoba}

\bkkigocards{\bkkigocard{$L$}{エル}{自己インダクタンス〔H〕。$N\Phi=Li$。第1巻の角運動量 $\vec L$、第3巻のラグランジアン $L$ とは別の量}{inductance}{電-8}}
{\bkkigocard{$M$}{エム}{相互インダクタンス〔H〕。$N_2\Phi_2=Mi_1$。電-7 の磁化 $\vec M$ とは別の量}{mutual inductance}{電-8}}
\bkkigocards{\bkkigocard{$\delta$}{デルタ}{表皮深さ〔m〕。$\delta=\sqrt{2/(\omega\mu\sigma)}$。第3巻の対数減衰率 $\delta$ とは別の量}{skin depth}{電-8}}{}

%% ======================== 例題 ========================
\Needspace{10\baselineskip}\bksection{reidai}{例題でたしかめよう}
\label{ph:電-8:ex1}
\begin{phreidai}[コイルを通る磁束が変わる]
100回巻きのコイルを通る磁束（1巻きあたり）が、\qty{0.10}{s}のあいだに、0から $\qty{1.0e-4}{Wb}$ まで一定の割合で増えました。コイルに生まれる誘導起電力の大きさを求めましょう。
\phbamen
コイルの中の磁束が増え続ける場面です。磁石をコイルに近づけるときのようすを、平均の変化で考えます。
\phmushi
コイルの抵抗と、誘導電流が作る磁場の影響（自己誘導）は無視します。磁束は0.10 s のあいだ、同じ割合で増えるとします。
\phzu
$N=100$、$\Delta\Phi=\qty{1.0e-4}{Wb}$、$\Delta t=\qty{0.10}{s}$。
\phriyu
コイルを通る磁束の変わる速さが与えられているので、ファラデーの法則（③）を使います。
\phtateru
\[ |\mathcal E|=N\frac{\Delta\Phi}{\Delta t} \]
\phtoku
\[ |\mathcal E|=100\times\frac{1.0\times10^{-4}}{0.10}=\qty{0.10}{V} \]
\phtashikame
単位：Wb/s＝V。\\
大きさ：乾電池（1.5 V）の約15分の1。やってみようのテスターの針がふれるくらいの電圧です。\\
極端な場合：磁束が増えるのに10倍の時間（1.0 s）がかかれば、電圧は10分の1になります。
\kotae \qty{0.10}{V}
\end{phreidai}
\begin{jibun}
200回巻きのコイルを通る磁束が、\qty{0.020}{s}のあいだに $\qty{3.0e-5}{Wb}$ 減りました。誘導起電力の大きさを求めましょう。増えるときと比べて、向きはどうなりますか。
\bkjans{$200\times3.0\times10^{-5}\div0.020=\qty{0.30}{V}$。向きは、増えるときと反対}
\end{jibun}

\label{ph:電-8:ex2}
\begin{phreidai}[レールの上を動く棒]
④の図のように、間かく\qty{0.50}{m}のレールに抵抗\qty{1.0}{\ohm}をつなぎ、向こう向きの\qty{0.10}{T}の磁場の中で、棒を右へ\qty{2.0}{m/s}で動かし続けます。誘導起電力、電流、棒を動かし続けるのに要る力、その仕事率を求め、ジュール熱と比べましょう。
\phbamen
磁場を横切って動く導体に電圧が生まれ、回路に電流が流れる場面です。
\phmushi
棒とレールの抵抗、棒とレールのこすれ、回路の電流が作る磁場（自己誘導）は無視します。
\phzu
④の図。$l=0.50$、$v=2.0$、$B=0.10$、$R=1.0$。
\phriyu
棒の起電力は $vBl$（④）、電流はオームの法則、力は電流が磁場から受ける力 $IlB$（電-6 ⑧）です。速さが一定なので、手の力は磁場の力とつり合います。
\phtateru
\[ \mathcal E=vBl,\qquad I=\frac{\mathcal E}{R},\qquad F=IlB,\qquad P=Fv \]
\phtoku
\[ \mathcal E=2.0\times0.10\times0.50=\qty{0.10}{V},\qquad I=\frac{0.10}{1.0}=\qty{0.10}{A} \]
\[ F=0.10\times0.50\times0.10=\qty{5.0e-3}{N},\qquad P=5.0\times10^{-3}\times2.0=\qty{0.010}{W} \]
ジュール熱は $RI^2=1.0\times0.10^2=\qty{0.010}{W}$ で、手の仕事率と等しくなります。
\phtashikame
単位：m/s×T×m＝V、A×m×T＝N。\\
大きさ：\qty{5}{mN}は、1円玉（\qty{1}{g}）の重さの約半分です。\\
極端な場合：抵抗をはずすと（回路を開く）、電流は流れず、力もいりません。電圧 0.10 V は、棒の両はしに生まれたままです。
\kotae $\mathcal E=\qty{0.10}{V}$、$I=\qty{0.10}{A}$、$F=\qty{5.0e-3}{N}$、仕事率 0.010 W（ジュール熱と等しい）
\end{phreidai}
\begin{jibun}
例題2 で、棒の速さを\qty{4.0}{m/s}にすると、起電力・電流・力・仕事率はどうなりますか。
\bkjans{$\qty{0.20}{V}$、$\qty{0.20}{A}$、$\qty{0.010}{N}$、$\qty{0.040}{W}$（仕事率は4倍）}
\end{jibun}

\label{ph:電-8:ex3}
\begin{phreidai}[発電機の電圧の山（3けたで書く）]
面積\qty{0.0100}{m^2}、100回巻きのコイルを、\qty{0.500}{T}の一様な磁場の中で、1秒に50回まわします。起電力の最大値を求めましょう（表計算の表\ref{ph:電-8:tab-sheet}と同じ数なので、3けたで書きます）。
\phbamen
一様な磁場の中でコイルを回し、交流を作る場面です（⑤）。
\phmushi
磁場はコイルの位置によらず一様とし、コイルの抵抗と自己誘導は無視します（電流を流さないときの電圧）。
\phzu
⑤の図。$N=100$、$S=0.0100$、$B=0.500$、$f=50$。
\phriyu
回るコイルの起電力は $NBS\omega\sin\omega t$（⑤）。山は $NBS\omega$ です。
\phtateru
\[ \mathcal E_{\max}=NBS\omega,\qquad \omega=2\pi f \]
\phtoku
\[ \omega=2\pi\times50\approx\qty{314}{rad/s},\qquad \mathcal E_{\max}=100\times0.500\times0.0100\times314.2\approx\qty{157}{V} \]
\phtashikame
単位：T×m$^2$×rad/s＝Wb/s＝V。\\
大きさ：家のコンセントの電圧の山（約\qty{141}{V}：電-10 ②）と同じくらいです。\\
極端な場合：回さなければ（$\omega=0$）、電圧は0です。
\kotae 約\qty{157}{V}
\end{phreidai}
\begin{jibun}
例題3 の発電機を、1秒に60回まわすと、山は何 V になりますか。
\bkjans{$\omega=2\pi\times60\approx377$ rad/s。山は約\qty{188}{V}（50 Hz のときの1.2倍）}
\end{jibun}

\label{ph:電-8:ex4}
\begin{phreidai}[飛行機の翼に生まれる電圧]
はば\qty{60}{m}の翼をもつ飛行機が、\qty{250}{m/s}で水平に飛んでいます。地球の磁場の上下の成分を $\qty{3.5e-5}{T}$ とします。翼の両はしのあいだに生まれる電圧を求めましょう。この電圧で、翼につないだ電球をともせますか。
\phbamen
地球の磁場の中を、導体（金属の翼）が横切る場面です。
\phmushi
翼を1本のまっすぐな導体の棒とみなします。地球の磁場の、水平の成分は、水平に動く翼の両はしには電圧を生まないので考えません。
\phzu
$l=\qty{60}{m}$、$v=\qty{250}{m/s}$、上下の成分 $B=\qty{3.5e-5}{T}$。棒・速度・磁場がたがいに直角。
\phriyu
磁場を横切る棒の起電力 $vBl$（④）を使います。
\phtateru
\[ \mathcal E=vBl \]
\phtoku
\[ \mathcal E=250\times3.5\times10^{-5}\times60\approx\qty{0.53}{V} \]
\phtashikame
単位：m/s×T×m＝V。\\
大きさ：乾電池の3分の1くらいです。\\
極端な場合：電球をつなぐ導線も、いっしょに飛んで同じ磁場を横切るので、導線にも同じ向きに同じ電圧が生まれます。閉じた回路を1周すると打ち消し合い、電流は流れません。電圧は生まれていても、電球はともせないのです（回路を通る磁束が変わらないから、と言っても同じです）。
\kotae 約\qty{0.53}{V}。電流は取り出せない
\end{phreidai}
\begin{jibun}
長さ\qty{1.5}{m}の新幹線の車軸が、時速\qty{300}{km}で走っています。地球の磁場の上下の成分を $\qty{3.5e-5}{T}$ として、車軸の両はしの電圧を求めましょう。
\bkjans{$300\div3.6\approx\qty{83}{m/s}$。電圧は $vBl$ で、約\qty{4.4e-3}{V}（約\qty{4.4}{mV}）}
\end{jibun}

\label{ph:電-8:ex5}
\begin{phreidai}[ソレノイドのインダクタンス]
長さ\qty{0.10}{m}、半径\qty{1.0}{cm}の筒に、エナメル線を1000回巻いたソレノイドの自己インダクタンスを求めましょう。
\phbamen
細長いコイルの形から、インダクタンスを見積もる場面です。
\phmushi
鉄心はない（空心）とし、はしの効果（はしで磁場が弱くなること）を無視します。長さは半径の10倍なので、長いソレノイドとみなします。
\phzu
$N=1000$、$l=\qty{0.10}{m}$、$S=\pi\times0.010^2\approx\qty{3.14e-4}{m^2}$。
\phriyu
長いソレノイドの $L=\mu_0N^2S/l$（⑥）が使えます。
\phtateru
\[ L=\mu_0\frac{N^2S}{l} \]
\phtoku
\[ L=4\pi\times10^{-7}\times\frac{1000^2\times3.14\times10^{-4}}{0.10}\approx\qty{3.9e-3}{H} \]
\phtashikame
単位：(N/A$^2$)×m$^2$÷m＝N·m/A$^2$＝J/A$^2$＝H（$U=\tfrac12LI^2$ から H＝J/A$^2$）。\\
大きさ：\qty{3.9}{mH}。数学の本の LC 回路（6-6）の\qty{10}{mH}と同じくらいの大きさです。\\
極端な場合：実際の、長さが半径の10倍のコイルでは、はしの効果で式より1割ほど小さくなります（専門の本の長岡係数）。
\kotae 約\qty{3.9}{mH}
\end{phreidai}
\begin{jibun}
例題5 のソレノイドの巻数を500回にすると、$L$ はいくらになりますか。また、1000回巻きのまま、比透磁率1000 の鉄心を入れると、いくらになりますか（鉄心は飽和しないとします）。
\bkjans{500回：約\qty{0.99}{mH}（4分の1）。鉄心：約\qty{3.9}{H}（1000倍）}
\end{jibun}

\label{ph:電-8:ex6}
\begin{phreidai}[リレーのコイルのスイッチを切る]
リレー（電磁石でスイッチを動かす部品）のコイルの自己インダクタンスを\qty{0.20}{H}、流れている電流を\qty{30}{mA}とします（数学 6-4 のリレーと同じ数）。スイッチを切って、電流が\qty{0.10}{ms}で0になったとすると、コイルに生まれる起電力の大きさは、平均でいくらですか。
\phbamen
コイルの電流を急に止める場面です。
\phmushi
電流は0.10 ms のあいだ、一定の割合で減るとします。コイルの抵抗は考えません。
\phzu
$L=\qty{0.20}{H}$、$\Delta I=\qty{30}{mA}=\qty{0.030}{A}$、$\Delta t=\qty{0.10}{ms}=\qty{1.0e-4}{s}$。
\phriyu
自己誘導の起電力 $L\,\Delta I/\Delta t$（⑥）です。
\phtateru
\[ |\mathcal E|=L\frac{\Delta I}{\Delta t} \]
\phtoku
\[ |\mathcal E|=0.20\times\frac{0.030}{1.0\times10^{-4}}=\qty{60}{V} \]
\phtashikame
単位：H×A/s＝V。\\
大きさ：リレーを動かす電源（12 V）の5倍の電圧が、スイッチの両はしにかかります。\\
極端な場合：切る時間が短いほど、電圧は大きくなります。スイッチのすき間に火花がとんで、接点がいたむことがあります。電-9 ④で、ダイオードでこれを防ぐ方法を学びます。
\kotae 約\qty{60}{V}
\end{phreidai}
\begin{jibun}
例題6 で、電流が $\qty{1.0}{\micro s}$ で0になると、起電力の大きさは平均でいくらですか。
\bkjans{$0.20\times0.030\div(1.0\times10^{-6})=\qty{6.0e3}{V}$（6000 V）}
\end{jibun}

\label{ph:電-8:ex7}
\begin{phreidai}[となりのコイルに生まれる電圧]
相互インダクタンスが\qty{0.010}{H}の2つのコイルがあります。コイル1の電流を、1秒に\qty{100}{A}の割合で増やすとき、コイル2に生まれる起電力の大きさを求めましょう。
\phbamen
一方のコイルの電流の変化で、もう一方のコイルに電圧が生まれる場面です（⑦）。
\phmushi
コイル2には電流が流れない（電圧計だけをつなぐ）とし、コイル2の電流がコイル1に及ぼす影響は考えません。
\phzu
$M=\qty{0.010}{H}$、$\dd i_1/\dd t=\qty{100}{A/s}$。
\phriyu
相互誘導の式 $\mathcal E_2=-M\,\dd i_1/\dd t$（⑦）です。
\phtateru
\[ |\mathcal E_2|=M\frac{\dd i_1}{\dd t} \]
\phtoku
\[ |\mathcal E_2|=0.010\times100=\qty{1.0}{V} \]
\phtashikame
単位：H×A/s＝V。\\
大きさ：電流を一定にすると0になり、変化を速くすると大きくなります。\\
極端な場合：$M=0$（2つのコイルの磁束がまったく通い合わない）なら、0です。
\kotae \qty{1.0}{V}
\end{phreidai}
\begin{jibun}
例題5 のソレノイド（1000回巻き・長さ\qty{0.10}{m}・半径\qty{1.0}{cm}）の上に、100回巻きのコイルを同じ長さに重ねて巻きました。相互インダクタンスを求め、ソレノイドの電流を1秒に\qty{50}{A}の割合で変えたときの、外側のコイルの起電力の大きさを求めましょう。
\bkjans{$M=\mu_0N_1N_2S/l$ で、約 $\qty{3.9e-4}{H}$。起電力は $M\times50$ で、約\qty{0.020}{V}}
\end{jibun}

\label{ph:電-8:ex8}
\begin{phreidai}[リレーのコイルにたまるエネルギー]
例題6 のリレーのコイル（\qty{0.20}{H}、\qty{30}{mA}）にたまっているエネルギーを求めましょう。
\phbamen
電流が流れているコイルの磁場に、エネルギーがたまっている場面です。
\phmushi
鉄心の飽和とヒステリシス（電-7 ⑧）は無視し、$L$ は一定とします。
\phzu
$L=\qty{0.20}{H}$、$I=\qty{0.030}{A}$。
\phriyu
コイルのエネルギーは $\tfrac12LI^2$（⑧）です。
\phtateru
\[ U=\frac12LI^2 \]
\phtoku
\[ U=\frac12\times0.20\times0.030^2=\qty{9.0e-5}{J} \]
\phtashikame
単位：H×A$^2$＝(V·s/A)×A$^2$＝V·A·s＝W·s＝J。\\
大きさ：とても小さなエネルギーですが、例題6 のように\qty{0.1}{ms}で放すと、仕事率は $9.0\times10^{-5}\div10^{-4}\approx\qty{1}{W}$ にもなり、小さな接点を傷めるのに十分です。\\
極端な場合：電流が2倍なら4倍です。
\kotae $\qty{9.0e-5}{J}$
\end{phreidai}
\begin{jibun}
例題5 のソレノイド（\qty{3.9}{mH}）に\qty{2.0}{A}を流します。（1）$\tfrac12LI^2$ でエネルギーを求めましょう。（2）中の磁場 $B=\mu_0nI$ と、磁場のエネルギー密度 $B^2/(2\mu_0)$ に中の体積をかけて求め、（1）と比べましょう。
\bkjans{（1）約 $\qty{7.9e-3}{J}$。（2）$B=\mu_0nI\approx\qty{0.025}{T}$、$u=B^2/(2\mu_0)\approx\qty{251}{J/m^3}$。体積 $\qty{3.14e-5}{m^3}$ をかけて、約 $\qty{7.9e-3}{J}$。（1）と同じ}
\end{jibun}

\label{ph:電-8:ex9}
\begin{phreidai}[電磁石が鉄を引く力]
電磁石の極の面積が $\qty{10}{cm^2}$ で、鉄の板とのすき間の磁場が\qty{1.0}{T}です。1つの極が鉄を引く力を求めましょう。
\phbamen
電磁石が鉄の板を引きつける場面です（⑩）。
\phmushi
すき間の磁場は一様で、すき間の外へのもれは無視します。鉄の中の磁場のエネルギーも無視します。
\phzu
$S=\qty{10}{cm^2}=\qty{1.0e-3}{m^2}$、$B=\qty{1.0}{T}$。
\phriyu
仮想変位で求めた $F=B^2S/(2\mu_0)$（⑩）です。
\phtateru
\[ F=\frac{B^2S}{2\mu_0} \]
\phtoku
\[ F=\frac{1.0^2\times1.0\times10^{-3}}{2\times4\pi\times10^{-7}}\approx\qty{4.0e2}{N} \]
\phtashikame
単位：T$^2$×m$^2$÷(N/A$^2$)。T＝N/(A·m) なので、$\mathrm{N^2/(A^2m^2)}\times\mathrm{m^2}\times\mathrm{A^2/N}=\mathrm{N}$。\\
大きさ：\qty{400}{N}は約\qty{40}{kg}の物の重さで、手のひらくらいの面積でも、強い力が出ます。極が2つ（U字形）なら2倍です。\\
極端な場合：すき間の磁場が0なら、力も0です。
\kotae 約 $\qty{4.0e2}{N}$
\end{phreidai}
\begin{jibun}
例題9 で、すき間の磁場を\qty{0.50}{T}にすると、力は何 N ですか。
\bkjans{$\tfrac14$ になって、約 $\qty{1.0e2}{N}$（99 N）}
\end{jibun}

\label{ph:電-8:ex10}
\begin{phreidai}[IH のなべの表皮深さ（はてな？）]
IH調理器が\qty{25}{kHz}で動いているとします。鉄（導電率 $\qty{1.0e7}{S/m}$、比透磁率100 とします）、アルミニウム（$\qty{3.8e7}{S/m}$）、銅（$\qty{6.0e7}{S/m}$）の表皮深さを求め、電流の流れる層の流れにくさ $1/(\sigma\delta)$ を、鉄と銅で比べましょう。
\phbamen
なべの底の金属に、高い周波数の磁場が入っていく場面です（⑫）。
\phmushi
なべの底は、平らで、表皮深さより十分厚いとします。鉄の比透磁率は一定とし、飽和は無視します。アルミと銅の比透磁率は1 とします。
\phzu
$f=\qty{25}{kHz}$、$\omega=2\pi f\approx\qty{1.57e5}{rad/s}$、$\mu=\mu_{\mathrm r}\mu_0$。
\phriyu
表皮深さは $\delta=\sqrt{2/(\omega\mu\sigma)}$（⑫）。層の流れにくさは、厚さ $\delta$ の板と同じと考えて $1/(\sigma\delta)$ に比例します。
\phtateru
\[ \delta=\sqrt{\frac{2}{\omega\mu_{\mathrm r}\mu_0\sigma}} \]
\phtoku
鉄では、$\omega\mu_{\mathrm r}\mu_0\sigma=1.57\times10^5\times100\times1.26\times10^{-6}\times1.0\times10^{7}\approx1.98\times10^{8}$ なので、
\[ \delta=\sqrt{\frac{2}{1.98\times10^{8}}}\approx\qty{1.0e-4}{m}=\qty{0.10}{mm} \]
同じように、アルミは約\qty{0.52}{mm}、銅は約\qty{0.41}{mm}です。流れにくさは、鉄が $1/(1.0\times10^7\times1.0\times10^{-4})\approx1.0\times10^{-3}$、銅が $1/(6.0\times10^7\times4.1\times10^{-4})\approx4.1\times10^{-5}$ で、鉄は銅の約24倍です。
\phtashikame
単位：$\sqrt{1/(\mathrm{s^{-1}}\times\mathrm{H/m}\times\mathrm{S/m})}$。H＝Ω·s、S＝1/Ω なので、$\sqrt{\mathrm{m^2}}=\mathrm{m}$。\\
大きさ：どれも\qty{1}{mm}より小さく、なべの底（数 mm）の、ほんの表面だけで熱が出ます。\\
極端な場合：周波数を下げていくと $\delta$ は大きくなり、50 Hz の銅では約\qty{9}{mm}です。
\kotae 鉄 約\qty{0.10}{mm}、アルミ 約\qty{0.52}{mm}、銅 約\qty{0.41}{mm}。鉄の流れにくさは銅の約24倍
\end{phreidai}
\begin{jibun}
アルミのなべを熱くするために、周波数を\qty{90}{kHz}にすると、アルミの表皮深さはいくらになりますか。25 kHz のときの何倍ですか。
\bkjans{約\qty{0.27}{mm}。$\sqrt{25/90}\approx0.53$ 倍}
\end{jibun}

\label{ph:電-8:ex11}
\begin{phreidai}[鉄心の板をうすくする]
変圧器の鉄心を、厚さ\qty{1.0}{mm}の板から\qty{0.35}{mm}の板に変えます。周波数と磁場の大きさが同じとき、鉄心全体の渦電流損は何倍になりますか。
\phbamen
鉄心の中で、渦電流によるむだな熱を減らす場面です（⑪）。
\phmushi
ヒステリシス損は考えません。板はどちらも表皮深さより十分うすく、鉄心全体の体積は同じとします。板の間の絶縁は完全とします。
\phzu
$d_1=\qty{1.0}{mm}$、$d_2=\qty{0.35}{mm}$。
\phriyu
うすい板の渦電流損は、1 m$^3$ あたり $d^2$ に比例します（⑪）。体積が同じなので、全体も $d^2$ に比例します。
\phtateru
\[ \frac{P_2}{P_1}=\Bigl(\frac{d_2}{d_1}\Bigr)^2 \]
\phtoku
\[ \frac{P_2}{P_1}=\Bigl(\frac{0.35}{1.0}\Bigr)^2\approx0.12 \]
\phtashikame
単位：同じ単位どうしの比なので、単位はありません。\\
大きさ：約8分の1。⑪のプログラムの 2.42 W/kg と 0.30 W/kg の比と同じです。\\
極端な場合：板を限りなくうすくすれば、渦電流損は0に近づきます（そのかわり、板の数と、絶縁の層が増えます）。
\kotae 約0.12倍（約8分の1）
\end{phreidai}
\begin{jibun}
（1）厚さ\qty{0.50}{mm}の板を\qty{0.35}{mm}にすると、渦電流損は何倍ですか。（2）同じ変圧器を、50 Hz の地域から60 Hz の地域に移すと、磁場の山が同じとき、渦電流損は何倍ですか。
\bkjans{（1）$(0.35/0.50)^2\approx0.49$ 倍。（2）$(60/50)^2\approx1.4$ 倍}
\end{jibun}

\label{ph:電-8:ex12}
\begin{phreidai}[IC カードのコイルと共振]
交通系 IC カードのコイルの自己インダクタンスを $\qty{2.0}{\micro H}$ とします。13.56 MHz で共振させるには、コンデンサーの静電容量をいくらにすればよいですか。
\phbamen
コイルとコンデンサーで、決まった周波数の磁場を効率よく受けとる場面です（⑬）。
\phmushi
コイルの抵抗と、チップの中の静電容量は無視します。
\phzu
$L=\qty{2.0e-6}{H}$、$f=\qty{13.56e6}{Hz}$。
\phriyu
コイルとコンデンサーの回路がゆれやすい周波数は $f=1/(2\pi\sqrt{LC})$（電-9 ⑦）です。
\phtateru
\[ f=\frac{1}{2\pi\sqrt{LC}}\qquad\text{より}\qquad C=\frac{1}{(2\pi f)^2L} \]
\phtoku
\[ C=\frac{1}{(2\pi\times13.56\times10^6)^2\times2.0\times10^{-6}}\approx\qty{6.9e-11}{F}=\qty{69}{pF} \]
\phtashikame
単位：$1/(\mathrm{s^{-2}}\times\mathrm{H})$。H×F＝s$^2$（数学 6-6）なので、F。\\
大きさ：pF（ピコファラド、$10^{-12}$ F）の大きさで、IC チップの中に作りこめる小さなコンデンサーです。\\
極端な場合：$L$ が大きいほど、要る $C$ は小さくなります。
\kotae 約\qty{69}{pF}
\end{phreidai}
\begin{jibun}
例題12 で、コイルを $\qty{1.0}{\micro H}$ にすると、要る静電容量はいくらですか。
\bkjans{約 $\qty{1.4e2}{pF}$（138 pF。2倍）}
\end{jibun}

%% ======================== まちがえやすいところ ========================
\Needspace{11\baselineskip}\bksection{tsumazuki}{まちがえやすいところをチェック}
\begin{tsumazuki}{強い磁石ほど、大きな電圧？}
\bkng{強い磁石をコイルの中に入れておけば、大きな電圧が生まれる。}
\bkok{電圧が生まれるのは、磁束が\textbf{変わっている}あいだだけ。止めておけば0。}
起電力を決めるのは、磁束の大きさではなく、変わる速さです（③）。
\end{tsumazuki}

\begin{tsumazuki}{誘導電流の磁場は、いつも磁石の磁場と反対向き？}
\bkng{誘導電流は、磁石の磁場を打ち消す向きの磁場を作る。}
\bkok{打ち消すのは磁場の\textbf{変化}。磁石を遠ざけるときは、磁石の磁場と同じ向きの磁場を作って、減るのを補う。}
レンツの法則は「変化をじゃまする」です（③の図\ref{ph:電-8:fig-lenz}の右）。
\end{tsumazuki}

\begin{tsumazuki}{磁束は、磁場×面積？}
\bkng{磁場 \qty{0.5}{T} の中に、面積 \qty{0.01}{m^2} のコイルがあるので、磁束はいつも $\qty{5e-3}{Wb}$。}
\bkok{$\Phi=BS\cos\theta$。面が磁場にそっていれば（$\theta=90^\circ$）、磁束は0。}
磁束は、面を\textbf{通りぬける}磁場の量です。発電機は、この角が変わることで電圧を生みます（⑤）。
\end{tsumazuki}

\begin{tsumazuki}{電圧があれば、電流が流れる？}
\bkng{飛行機の翼には 0.5 V の電圧があるので、導線でつなげば電流がとれる。}
\bkok{つないだ導線にも同じ電圧が生まれ、回路を1周すると打ち消し合う。回路を通る磁束が変わらなければ、電流は流れない。}
例題4 を見直しましょう。
\end{tsumazuki}

\begin{tsumazuki}{インダクタンスは、電流で変わる？}
\bkng{電流を2倍にすると、磁束も2倍になるので、$L$ も2倍になる。}
\bkok{$L=N\Phi/i$ は比なので、電流によらない（鉄心がないとき）。$L$ はコイルの形と巻数で決まる。}
ただし、鉄心が飽和すると、電流を増やしても磁束が増えにくくなり、$L$ は小さくなります（⑥のモデル化の札）。
\end{tsumazuki}

\begin{column}{ファラデーの鉄の輪}
ファラデーは、ロンドンの王立研究所の研究者でした。1831年の夏、やわらかい鉄の輪の片側にコイルを巻いて電池につなぎ、反対側のコイルを検流計につなぎました。電池をつないだ瞬間と、はずした瞬間にだけ、針がふれたのです。この輪は、今も王立研究所に残っています。2つのコイルを鉄心に巻いたこの形は、今の変圧器（電-12 ⑤）とまったく同じです。ファラデーは、このあと、磁石をコイルに出し入れする実験や、銅の円板を磁石の極のあいだで回して電流を取り出し続ける実験を、次々に行いました。
\end{column}

%% ======================== 工学ではこう使う ========================
\Needspace{18\baselineskip}\bksection{phgear}{工学ではこう使う}
\begin{phkougaku}[電気をつくる・送る・使う、どこにも電磁誘導]
\textbf{つくる}：発電所の発電機は、回る電磁石の周りのコイルに交流を生ませます（⑤）。電気自動車や電車は、ブレーキをかけるとき、モーターを発電機として使い、車の運動エネルギーを電気にもどします（回生ブレーキ：電-12）。

\textbf{送る}：変圧器は、相互誘導で電圧を上げ下げします（⑦）。鉄心はうすい電磁鋼板を重ねて、渦電流損を小さくしています（⑪）。

\textbf{使う・調べる}：IH調理器（⑫）、IC カードとワイヤレス充電（⑬）のほか、金属の管や板の中の傷を、コイルで渦電流を流して見つける\textbf{渦流探傷}、空港の金属探知機、電流を電線をはさんで測るクランプメーター（交流用。相互誘導）も、電磁誘導を使っています。

\begin{phketa}[電磁誘導の数の目安]
\phketaline{IH調理器の周波数}{20.05〜100 kHz（電波法施行規則）}
\phketaline{IC カード}{13.56 MHz（コイルは \si{\micro H}、コンデンサーは pF）}
\phketaline{銅の表皮深さ}{50 Hz で約9 mm、25 kHz で約0.4 mm}
\phketaline{磁場が引く力}{1 T のとき、1 cm$^2$ あたり約40 N}
\phboxhead{答えが桁で合っているか}
誘導起電力は「巻数×磁束÷時間」です。手で動かす磁石とコイルなら、磁束は $10^{-5}$〜$10^{-4}$ Wb、時間は0.1 s くらいなので、巻数100回で0.01〜0.1 V です。何 kV のような答えになったら、時間か面積の単位（ms、cm$^2$）の直し忘れを疑いましょう。
\end{phketa}
\end{phkougaku}
\phstation{電気機器,電気回路,非破壊検査,電磁界解析}{誘導起電力、鎖交磁束、自己インダクタンス、相互インダクタンス、渦電流損、表皮効果、逆起電力}

%% ======================== 章末問題 ========================
\bkendproblems
\bkproblemgroup{基本}{3}
\bkproblem 50回巻きのコイルを通る磁束が、$\qty{4.0e-3}{Wb}$ から\qty{0.20}{s}で0まで、一定の割合で減りました。誘導起電力の大きさを求めましょう。
\bkpans{$50\times4.0\times10^{-3}\div0.20=\qty{1.0}{V}$}
\bkproblem 間かく\qty{0.30}{m}のレールの上の棒を、\qty{0.20}{T}の磁場に直角に\qty{5.0}{m/s}で動かします。回路の抵抗は\qty{0.50}{\ohm}です。起電力・電流・棒を動かし続ける力・その仕事率を求めましょう。
\bkpans{$\mathcal E=\qty{0.30}{V}$、$I=\qty{0.60}{A}$、$F=IlB=\qty{0.036}{N}$、仕事率 $\qty{0.18}{W}$（$RI^2$ と等しい）}
\bkproblem 200回巻き、面積\qty{0.020}{m^2}のコイルを、\qty{0.25}{T}の磁場の中で1秒に50回まわします。起電力の山を求めましょう。
\bkpans{$NBS\omega$ に入れて、約 $\qty{3.1e2}{V}$（314 V）}
\bkproblem 長さ\qty{0.20}{m}、断面積 $\qty{2.0}{cm^2}$、400回巻きの空心のソレノイドがあります。（1）自己インダクタンスを求めましょう。（2）\qty{3.0}{A}を流したとき、たまるエネルギーを求めましょう。
\bkpans{（1）$\mu_0N^2S/l$ に入れて、約 $\qty{2.0e-4}{H}$。（2）$\tfrac12LI^2$ で、約 $\qty{9.0e-4}{J}$}
\bkproblem 相互インダクタンスが\qty{0.050}{H}の2つのコイルがあります。一方の電流を、\qty{0.010}{s}のあいだに\qty{2.0}{A}だけ一定の割合で増やすと、もう一方に生まれる起電力の大きさはいくらですか。
\bkpans{$0.050\times2.0\div0.010=\qty{10}{V}$}
\bkproblem たてに立てた2本のレール（間かく\qty{0.20}{m}）の上はしを\qty{0.10}{\ohm}の抵抗でつなぎ、レールにそって、質量\qty{0.010}{kg}の金属の棒を水平にして静かに放します。レールの面に垂直に\qty{0.50}{T}の磁場があります。棒の速さは、やがて一定になります。その速さを求めましょう（こすれは無視、$g=\qty{9.8}{m/s^2}$）。
\bkpans{一定の速さでは、重さと磁場の力がつり合う：$mg=IlB=(vBl/R)\,lB$。$v=mgR/(B^2l^2)=\qty{0.98}{m/s}$}
\bkproblemgroup{チャレンジ}{5}
\bkproblem モーターの鉄心を、厚さ\qty{0.50}{mm}の板から\qty{0.25}{mm}の板に変え、同時に、回す周波数を50 Hz から400 Hz に上げます。磁場の山が同じとき、1 m$^3$ あたりの渦電流損は何倍になりますか。板は表皮深さより十分うすいとします。
\bkpans{$(0.25/0.50)^2\times(400/50)^2=0.25\times64=16$ 倍。周波数を上げると、板をうすくしても損失が増える}
\bkproblem 半径\qty{1.0}{mm}の銅線（導電率 $\qty{6.0e7}{S/m}$）に、1 MHz の交流を流します。（1）表皮深さを求めましょう。（2）電流が表面から深さ $\delta$ の層だけを一様に流れると考えると、電線の抵抗は直流のときの約何倍になりますか（層の断面積を $2\pi r\delta$ で近似する）。
\bkpans{（1）$\delta=\sqrt{2/(\omega\mu_0\sigma)}$ で、約\qty{0.065}{mm}。（2）断面積の比 $\pi r^2/(2\pi r\delta)=r/(2\delta)$ で、約7.7倍}
\bkproblem U字形の電磁石の2つの極（面積はどちらも $\qty{4.0}{cm^2}$）に、鉄の板が吸いついています。（1）すき間の磁場が\qty{0.80}{T}のとき、引く力の合計を求めましょう。（2）質量\qty{50}{kg}の鉄の板を持ち上げるには、すき間の磁場を少なくともいくらにする必要がありますか（$g=\qty{9.8}{m/s^2}$）。
\bkpans{（1）1極は $B^2S/(2\mu_0)$ で約\qty{102}{N}、2極で約 $\qty{2.0e2}{N}$。（2）$2\times B^2S/(2\mu_0)=50\times9.8$ より、$B\approx\qty{1.24}{T}$。持ち上げるには、これより大きく（約\qty{1.3}{T}）する}

\begin{furikaeri}
\begin{checklist}
\item 磁石をコイルに出し入れすると電流が流れ、動かす向きと速さで、電流の向きと大きさが変わることを、実験で確かめられる。\bkfurref{ph:電-8:sec1}
\item ファラデーの法則とレンツの法則で、動く棒や発電機に生まれる電圧（誘導起電力）の大きさと向きを求められる。\bkfurref{ph:電-8:sec3}
\item 自己インダクタンス・相互インダクタンスを求め、コイルにたまるエネルギーを計算できる。\bkfurref{ph:電-8:sec6}
\item ファラデーの法則の微分形を読み、磁場のエネルギーから電磁石が鉄を引く力を求められる。\bkfurref{ph:電-8:sec9}
\item 渦電流と表皮効果で、IH調理器や、うすい板を重ねた鉄心のしくみを説明できる。\bkfurref{ph:電-8:sec12}
\end{checklist}
\end{furikaeri}
\bkfukidashi{mina}{IH のなべは、底に流れるうずの電流で、自分で熱くなっていたんだね。鉄のなべがよく熱くなるのは、電流がうすい表面に集まって、流れにくくなるからなんて、おもしろい！}
